% Source: preamble
\documentclass[11pt,letterpaper]{article}
\usepackage[margin=1.05in]{geometry}
\usepackage[T1]{fontenc}
\usepackage{lmodern,microtype}
\usepackage{amsmath,amssymb,amsthm,mathtools}
\usepackage{xcolor}
% Set \showpolishfalse below to obtain a clean copy.
\newif\ifshowpolish
\showpolishfalse
\newcommand{\polish}[1]{\ifshowpolish{\color{magenta}%
\hypersetup{linkcolor=magenta,citecolor=magenta,urlcolor=magenta}#1}\else#1\fi}
% Revision of 25 October 2026, addressing the major suggestions of the
% referee-style evaluation. All text that differs from the draft of
% 21 October 2026 can be shown in blue; deleted text is not shown. The
% marks are switched off; set \showrevtrue below to show them. The
% earlier magenta marks are switched off above.
\newif\ifshowrev
\showrevfalse
\definecolor{revblue}{RGB}{0,70,200}
\newcommand{\rev}[1]{\ifshowrev{\color{revblue}%
\hypersetup{linkcolor=revblue,citecolor=revblue,urlcolor=revblue}#1}\else#1\fi}
\newcommand{\rs}[1]{\ifshowrev{\color{revblue}#1}\else#1\fi}
% Notes on removed material; hidden by default (\showrevnotestrue shows them).
\newif\ifshowrevnotes
\showrevnotesfalse
\newcommand{\revnote}[1]{\ifshowrevnotes{\color{revblue}\footnotesize[#1]}\fi}
\newenvironment{revblock}{\ifshowrev\color{revblue}%
\hypersetup{linkcolor=revblue,citecolor=revblue,urlcolor=revblue}\fi}{}
\newcommand{\Ac}{\rs{\mathsf A}}
% Pass of 25 October 2026 against main (commit 1635fc3): every change of
% this pass can be shown in blue; the marks are off (\showpasstrue shows
% them).
\newif\ifshowpass
\showpassfalse
\newcommand{\pass}[1]{\ifshowpass{\color{revblue}%
\hypersetup{linkcolor=revblue,citecolor=revblue,urlcolor=revblue}#1}\else#1\fi}
\newcommand{\ps}[1]{\ifshowpass{\color{revblue}#1}\else#1\fi}
% Removal of definitions that the appendices repeated from the main text,
% and the restructuring for readability that followed (25 October 2026,
% against commit ab809b7): new text is shown in blue, deleted text is not
% shown, and the appendices were reordered without other changes. The
% marks are off; \showtrimtrue shows them.
\newif\ifshowtrim
\showtrimfalse
\newcommand{\trim}[1]{\ifshowtrim{\color{revblue}%
\hypersetup{linkcolor=revblue,citecolor=revblue,urlcolor=revblue}#1}\else#1\fi}
\newcommand{\trs}[1]{\ifshowtrim{\color{revblue}#1}\else#1\fi}
% Define-before-use pass (25 October 2026, against commit 6cc00d3): text
% that was added or moved so that every notion is defined at or before
% its first use is shown in blue when \showlintrue is set; the marks are
% off.
\newif\ifshowlin
\showlinfalse
\newcommand{\lin}[1]{\ifshowlin{\color{revblue}%
\hypersetup{linkcolor=revblue,citecolor=revblue,urlcolor=revblue}#1}\else#1\fi}
\newcommand{\lns}[1]{\ifshowlin{\color{revblue}#1}\else#1\fi}
% Precision pass after the rerun of the revision skill (25 October 2026,
% against commit 43d2670): the changes of the pass (the regime of the
% optimality claim, the description of the remainder, of the tube width and
% of the two ranges of \lambda in Section 2, the exponents, the citation of
% McFadden and the wording about Gao and Zhang in Subsection 2.6, the titles,
% the citations, the footnote on the AI collaborators, the shortened Lean
% section, the shortened abstract, and the renamed symbols) are shown in
% blue when \showskilltrue is set; the marks are off.
\newif\ifshowskill
\showskillfalse
\newcommand{\skl}[1]{\ifshowskill{\color{revblue}%
\hypersetup{linkcolor=revblue,citecolor=revblue,urlcolor=revblue}#1}\else#1\fi}
\newcommand{\sks}[1]{\ifshowskill{\color{revblue}#1}\else#1\fi}
\newenvironment{skillblock}{\ifshowskill\color{revblue}%
\hypersetup{linkcolor=revblue,citecolor=revblue,urlcolor=revblue}\fi}{}
% Readable geometry in the proof of Lemma C.5 (25 October 2026, against
% commit 87fd7a7): a reminder of the terms from the geometry of polyhedra, a
% rewritten first paragraph of the proof and a rewritten Step 1 are shown in
% blue when \showgeotrue is set; the marks are off.
\newif\ifshowgeo
\showgeofalse
\newenvironment{geoblock}{\ifshowgeo\color{revblue}%
\hypersetup{linkcolor=revblue,citecolor=revblue,urlcolor=revblue}\fi}{}
\definecolor{paperlink}{RGB}{0,0,114}
\usepackage[colorlinks=true,hypertexnames=false,linktoc=all,linkcolor=paperlink,citecolor=paperlink,urlcolor=paperlink]{hyperref}
\usepackage{fancyhdr}
\pagestyle{fancy}
\fancyhf{}
\fancyhead[L]{\small Gaussian approximation over rectangles}
\fancyhead[R]{\small\thepage}
\setlength{\headheight}{14pt}
\hypersetup{pdftitle={High-Dimensional Gaussian Approximations over Rectangles without Covariance Restrictions},pdfsubject={Bounded independent summands; rates optimal up to logarithmic factors},pdfauthor={Victor Chernozhukov, Denis Chetverikov, Kengo Kato, Yuta Koike}}
\setlength{\parindent}{1.25em}
\setlength{\parskip}{0pt}
\emergencystretch=2em
\allowdisplaybreaks[2]
\raggedbottom
\clubpenalty=10000
\widowpenalty=10000
\newtheorem{theorem}{Theorem}
\newtheorem{proposition}{Proposition}
\newtheorem{lemma}{Lemma}
\newtheorem{corollary}{Corollary}
\theoremstyle{remark}
\newtheorem{remark}{Remark}
\newcommand{\E}{\mathbb E}
\newcommand{\Pp}{\mathbb P}
\newcommand{\R}{\mathbb R}
\newcommand{\cR}{\mathcal R}
\newcommand{\cF}{\mathcal F}
\newcommand{\cP}{\mathcal P}
\newcommand{\cD}{\mathcal D}
\newcommand{\Cov}{\operatorname{Cov}}
\newcommand{\Var}{\operatorname{Var}}
\newcommand{\tr}{\operatorname{tr}}
\newcommand{\op}{\mathrm{op}}
\newcommand{\norm}[1]{\left\lVert#1\right\rVert}
\newcommand{\rh}{\rho_{\mathcal R}}
\newcommand{\ind}{\mathbf 1}
\newcommand{\dd}{\,d}
\newcommand{\Hh}{H_2}
\newcommand{\surf}{\mathcal J}

\newcommand{\logp}{\log_+}


% Source: front
\begin{document}
\begin{center}
{\Large High-Dimensional Gaussian Approximations\\[0.15em]
over Rectangles without Covariance Restrictions}\\[0.85em]
{\normalsize Victor Chernozhukov, Denis Chetverikov, Kengo Kato, Yuta
Koike\renewcommand{\thefootnote}{\fnsymbol{footnote}}\footnote{\skl{The paper
was developed through collaboration with the AI systems GPT-6 Astra,
GPT-6 Sol, and Claude Opus.}}}\\[0.65em]
{\normalsize 21 October 2026}
\end{center}
\setcounter{footnote}{0}
\thispagestyle{plain}

\begin{abstract}
\begin{skillblock}
We study the Gaussian approximation of sums of independent bounded random
vectors over rectangles, with no assumption on the covariance matrix,
which may be singular. Let \(X_1,\ldots,X_n\) be independent centered
random vectors in \(\R^p\) whose coordinates \(X_{ij}\) are bounded in
absolute value by a constant~\(B\), and suppose that each coordinate of
\(W_n=n^{-1/2}(X_1+\cdots+X_n)\) has unit variance. Let \(Z_n\) be a
centered Gaussian vector with the covariance matrix of \(W_n\), and put
\(L=\log(ep)\) and \(\log_+x=\max\{0,\log x\}\). The largest difference
between \(\Pp(W_n\in A)\) and \(\Pp(Z_n\in A)\) over rectangles \(A\) is
at most
\[
 \frac{C_BL^{3/2}}{\sqrt n}\sqrt{1+\log_+\sqrt{n/L}},
\]
where \(C_B\) depends only on~\(B\). If, for a fixed \(c_0>0\), no two
coordinates have correlation larger than \(1-c_0\) in absolute value, the
bound improves slightly, by a logarithmic factor. When \(p\) grows
polynomially in~\(n\), the two bounds are of order \(n^{-1/2}(\log n)^2\) and
\(n^{-1/2}(\log n)^{3/2}\), respectively. A further improvement holds if the
third moments cancel in aggregate, that is, if
\(\sum_i\E[X_{ij}X_{ik}X_{i\ell}]=0\) for all coordinates \(j,k,\ell\).
The power
\(n^{-1/2}\) cannot be improved even when \(p=1\), and a two-point lower
bound shows that the factor \(L^{3/2}\) cannot be improved either, so the
second rate is then optimal and the first is optimal up to the factor
\((\log n)^{1/2}\).
\end{skillblock}
\end{abstract}

\clearpage
\tableofcontents
\clearpage

\section{\texorpdfstring{\sks{Introduction and main results}}{Introduction and main results}}\label{sec:results}
Gaussian approximation over rectangles provides a common foundation for
approximating maxima and constructing simultaneous confidence regions.
\rev{In these problems the coordinates may satisfy exact linear relations,
so that the covariance matrix is singular although every coordinate has
positive variance, and nearly exact relations produce arbitrarily small
eigenvalues. We give bounds whose constants do not depend on the
eigenvalues.}

Let \(\mathcal R\) denote the class of axis-aligned rectangles in
\(\R^p\). We measure the approximation error by
\[
 \rh(U,V)=\sup_{A\in\mathcal R}
       |\Pp(U\in A)-\Pp(V\in A)|.
\]
For example, taking \(A=(-\infty,t]^p\) shows that this distance
controls the error in approximating the distribution of the largest
coordinate.

Consider a sum of independent centered random vectors,
\begin{equation}
 \begin{gathered}
 W=\sum_{i=1}^n\xi_i,\qquad
 \E\xi_i=0,\qquad \norm{\xi_i}_\infty\le\delta\quad\text{a.s.},\\
 \Sigma=\Cov(W),\qquad\Sigma_{jj}=1\ (j\le p),
 \qquad Z\sim N(0,\Sigma).
 \end{gathered}
 \label{eq:setup}
\end{equation}
The summands need not have the same distribution, and coordinates within
a summand may be dependent. \rev{The unit marginal variances fix the scale
of each coordinate; they can be obtained by coordinatewise rescaling
whenever all marginal variances are positive, and they place no
restriction on the rank of \(\Sigma\).}

The parameter \(\delta\) bounds the contribution of any one summand
to any coordinate of \(W\). For a normalized sample sum with
observations bounded by \(B\), it equals \(B/\sqrt n\). \rev{A uniform
bound of this kind is called an \emph{envelope}.} To state the
bound compactly, put
\[
 L=\log(ep),\qquad
 D_L=1+\log_+\frac1{\delta\sqrt L},\qquad
 \log_+x=\max\{0,\log x\}.
\]
Here \(L\) records the dependence on dimension, while \(D_L\)
contributes an additional logarithmic factor. Constants denoted by
\(C\) are absolute unless a subscript records their dependence.

\rev{Our first bound depends only on the summand envelope and the
dimension, so it holds uniformly as covariance eigenvalues approach or
reach zero.}

\begin{theorem}[Unrestricted covariance]\label{thm:unrestricted}
Under~\eqref{eq:setup},
\begin{equation}
 \rh(W,Z)\le C\min\{1,\delta L^{3/2}\sqrt{D_L}\}.
 \label{eq:unrestricted}
\end{equation}
\end{theorem}

\rev{Theorem~\ref{thm:unrestricted} allows any covariance rank, perfect
correlations, and singular principal submatrices, and for normalized
observations it gives the following dependence on the sample size.}

\begin{corollary}[Normalized observations]\label{cor:observations}
Let \(X_1,\ldots,X_n\in\R^p\) be independent and centered. Suppose
\(\norm{X_i}_\infty\le B\) almost surely and
\(n^{-1}\sum_i\E X_{ij}^2=1\) for every \(j\). Then, with
\(W_n=n^{-1/2}\sum_iX_i\) and \(Z_n\sim N(0,\Cov(W_n))\),
\begin{equation}
 \rh(W_n,Z_n)\le C\min\!\left\{1,
 \frac{B L^{3/2}}{\sqrt n}
 \sqrt{1+\log_+\frac{\sqrt n}{B\sqrt L}}\right\}.
 \label{eq:observations-unrestricted}
\end{equation}
\end{corollary}

\rev{Corollary~\ref{cor:observations} is Theorem~\ref{thm:unrestricted}
with \(\xi_i=X_i/\sqrt n\)\pass{, and for} fixed \(B\) and polynomial
dimension \(p\asymp n^a\), where \(a>0\) is fixed, \pass{its} bound has
order \(n^{-1/2}(\log n)^2\), so the usual power \(n^{-1/2}\) of the
sample size is kept up to logarithmic factors, however small the
covariance eigenvalues are.} \pass{The factor \(BL^{3/2}\) cannot be
improved, because Chernozhukov, Chetverikov, and
Koike~\cite[Proposition~2.1]{CCK} construct i.i.d.\ observations with
independent coordinates whose error is at least of order
\(BL^{3/2}/\sqrt n\); Theorem~\ref{opt:skew-lower} gives a version of
their example with the envelope attained. The bound
\eqref{eq:observations-unrestricted} is thus optimal up to the factor
\(\sqrt{1+\log_+(\sqrt n/(B\sqrt L))}\), which is of order
\((\log n)^{1/2}\) in polynomial dimension.} \skl{Whether this factor is
needed is open.}

For unequal marginal variances, define
\[
 \sigma_j^2=n^{-1}\sum_i\E X_{ij}^2,
 \qquad
 \underline\sigma=\min_j\sigma_j>0.
\]
Apply the result to \(X_{ij}/\sigma_j\). If
\(\max_i\norm{X_i}_\infty\le K_n\), the standardized summands have
envelope \(\delta=K_n/(\underline\sigma\sqrt n)\).
Positive diagonal scaling preserves rectangles and their approximation
distance, so the resulting bound also applies in the original coordinates.

\subsection{Correlation gaps and third-moment cancellation}
\polish{A better bound is available when no two distinct coordinates are nearly
equal or opposite. To measure this separation, define}
\[
 c=1-\max_{j\ne k}|\Sigma_{jk}|,
 \qquad c=1\ \text{when }p=1,
\]
and put
\begin{equation}
 \kappa=\begin{cases}\min\{L,c^{-1}\},&c>0,\\L,&c=0,\end{cases}
 \qquad E_\kappa=1+\log_+\frac{\sqrt\kappa}{\delta L}.
 \label{eq:parameters}
\end{equation}
A positive gap excludes equal or opposite coordinates; it still permits
singular covariance matrices.

\begin{theorem}[Gap-sensitive bound]\label{thm:optimized}
Under~\eqref{eq:setup},
\begin{equation}
 \rh(W,Z)\le C\min\!\left\{1,
 \delta[L^{3/2}+L\sqrt{\kappa E_\kappa}]\right\}.
 \tag{R*}\label{eq:optimized}
\end{equation}
\end{theorem}

\rev{Since \(\kappa\le L\) and \(E_\kappa\le D_L\),
Theorem~\ref{thm:optimized} implies Theorem~\ref{thm:unrestricted}\pass{;
the} term \(\delta L^{3/2}\) \pass{in~\eqref{eq:optimized}} comes from the
\lin{third moments of the summands},
and it disappears when the full third-moment tensor cancels.}

\begin{theorem}[Cancellation of third moments]\label{thm:zero-skew}
Under~\eqref{eq:setup}, suppose
\begin{equation}
 \sum_i\E\xi_i^{\otimes3}=0,
 \quad\text{equivalently,}\quad
 \sum_i\E[\xi_{ij}\xi_{ik}\xi_{i\ell}]=0
 \quad(j,k,\ell\le p).
 \label{eq:zero-third}
\end{equation}
Then
\begin{equation}
 \rh(W,Z)\le C\min\{1,\delta L\sqrt{\kappa E_\kappa}\}.
 \tag{ZS}\label{eq:zero-skew}
\end{equation}
\end{theorem}

\rev{Condition~\eqref{eq:zero-third} holds when each summand is centrally
symmetric, that is, when \(\xi_i\) and \(-\xi_i\) have the same law; it
involves the mixed third moments, so vanishing marginal
third moments do not suffice. The condition allows cancellation between
different summands and is preserved by deterministic linear maps.}

For a fixed standardized envelope and \(c\ge c_0>0\), the sample-size
bounds become
\begin{align}
 \rh(W_n,Z_n)&\le C_{B,c_0}\min\!\left\{1,
 \frac{L^{3/2}}{\sqrt n}
 +\frac L{\sqrt n}\sqrt{1+\log_+\frac{\sqrt n}{L}}\right\},
 \label{eq:observations-fixed}\\
 \rh(W_n,Z_n)&\le C_{B,c_0}\min\!\left\{1,
 \frac L{\sqrt n}\sqrt{1+\log_+\frac{\sqrt n}{L}}\right\}
 \quad\text{under~\eqref{eq:zero-third}}.
 \label{eq:zero-fixed}
\end{align}
\pass{In polynomial dimension the bound~\eqref{eq:observations-fixed}
has order \(n^{-1/2}(\log n)^{3/2}\), and it is optimal there, since the
lower-bound examples of~\cite[Proposition~2.1]{CCK} and
Theorem~\ref{opt:skew-lower} have independent coordinates and hence
correlation gap one.} These bounds give uniform convergence at a fixed
positive gap when \(\log p=o(n^{1/3})\) in general and \(\log p=o(\sqrt n)\)
under cancellation. Theorem~\ref{opt:consistency} shows that these
conditions are sharp over fixed-envelope classes, with envelope at least
two for the first assertion and at least one for the second. \rev{In
the second case the sharpness concerns only the dimension threshold, and
below the threshold the error rate in~\eqref{eq:zero-fixed} is not known
to be sharp.}

\subsection{Sub-exponential coordinates}\label{sec:subexponential}
\rev{The bounded theorems extend to sub-exponential coordinates if each
coordinate is clipped, that is, truncated to an interval \([-T,T]\), at a
level \(T\) that depends only on \(\log n\), and the discarded sum is
controlled separately. We measure tails by the Orlicz
norm}
\[
 \|V\|_{\psi_1}=\inf\{t>0:\E e^{|V|/t}\le2\}.
\]

\begin{corollary}[Sub-exponential observations]\label{cor:subexponential}
Let \(X_1,\ldots,X_n\in\R^p\) be independent and centered, with
\(n^{-1}\sum_i\E X_{ij}^2=1\) and
\(\max_{i,j}\|X_{ij}\|_{\psi_1}\le K\), where \(K\ge1\).
For \(W_n=n^{-1/2}\sum_iX_i\) and a centered Gaussian \(Z_n\)
with the same covariance, put
\[
 \ell_n=\log(en),\qquad b_n=\frac{K\ell_n}{\sqrt n},\qquad
 D_n=1+\log_+\frac1{b_n\sqrt L}.
\]
Then
\begin{equation}
 \rh(W_n,Z_n)\le C\min\{1,b_nL^{3/2}\sqrt{D_n}\}.
 \label{eq:subexponential}
\end{equation}
If the correlation gap of \(W_n\) is at least a fixed \(c_0>0\), then
\begin{equation}
 \rh(W_n,Z_n)\le C_{c_0}\min\!\left\{1,
 b_n\left[L^{3/2}+L\sqrt{1+\log_+\frac1{b_nL}}\right]\right\}.
 \label{eq:subexponential-gap}
\end{equation}
No local determinant or positive-eigenvalue lower bound is needed for
\eqref{eq:subexponential}. Unequal marginal variances are handled by
imposing the tail condition on \(X_{ij}/\sigma_j\), as in the
standardization following Corollary~\ref{cor:observations}.
\end{corollary}

For fixed \(K\) and \(p\asymp n^a\), \pass{where} \(a>0\) \pass{is
fixed}, these rates are
\(n^{-1/2}(\log n)^3\) and
\(n^{-1/2}(\log n)^{5/2}\), respectively\rev{, so the power of the
sample size is again optimal up to logarithmic factors}. In particular,
\eqref{eq:subexponential} guarantees convergence for fixed \(K\) when
\(\log p=o(n^{1/3}/\log(en))\). \rev{We do not claim that these
dimension ranges are sharp for the sub-exponential class.}
Appendix~\ref{app:subexponential} controls the discarded sum directly,
which allows the \lin{truncation level} to depend only on \(\log n\), and restores
covariance through a joint Gaussian coupling.

\subsection{Relation to recent work}\label{sec:literature}
High-dimensional Gaussian approximation for maxima and rectangles was
developed by Chernozhukov, Chetverikov, and Kato~\cite{CCK2013,CCK2017}.
Without covariance restrictions beyond positive marginal variances,
Chernozhukov, Chetverikov, Kato, and Koike~\cite[Theorem~2.1]{CCKK2022}
obtained an \(n^{-1/4}\) rate, up to logarithms, under coordinatewise
sub-exponential tails and an averaged fourth-moment condition.
Gao and Zhang~\cite[Theorem~1 and Section~5]{GaoZhang} recently improved
the sample-size power to \(n^{-1/3}\). Under the tail assumption of
Corollary~\ref{cor:subexponential}, their rectangle bound, in our
unit-variance notation, is
\begin{equation*}
 C\min\left\{1,
 \left(\frac{K^2}{n}\right)^{1/3}\{\log(2np)\}^{7/3}
 +\frac K{\sqrt n}\{\log(2np)\}^{5/2}\right\}.
\end{equation*}
No additional fourth-moment condition is needed for this bound.
For fixed \(K\) and \(p\asymp n^a\), with fixed \(a>0\), it has
order \(n^{-1/3}(\log n)^{7/3}\). In this dimension regime,
Corollaries~\ref{cor:observations} and~\ref{cor:subexponential}, with
fixed \(B\) or \(K\), give the sharp sample-size power
\(n^{-1/2}\), with logarithmic factors \((\log n)^2\) for bounded
observations and \((\log n)^3\) for sub-exponential coordinates.
These comparisons concern fixed envelopes or tail scales; the full
bounds retain their different dependence on these parameters.

A second line of work obtains rates close to \(n^{-1/2}\) under
covariance nondegeneracy. For bounded sums with positive definite
\(\Sigma\), Aronow and Lopatto~\cite[Theorem~1.1]{AL} remove the
\(\log(en)\) factor from the bound of Chernozhukov, Chetverikov,
and Koike~\cite[Corollary~2.1]{CCK}, and prove
\begin{equation*}
 \rh(W_n,Z_n)\le C\min\left\{1,
 \frac{BL^{3/2}}{\lambda_{\min}(\Sigma)\sqrt n}\right\}.
\end{equation*}
A fixed positive minimum eigenvalue therefore gives the sharp
\(BL^{3/2}/\sqrt n\) scale in the regimes of matching lower
bounds\pass{~\cite[Proposition~2.1]{CCK}}.
\rev{Our unrestricted bound does not involve \(\lambda_{\min}(\Sigma)\),
and it exceeds this scale only by the factor \(\sqrt{D_L}\).}

\rev{Other results replace full nondegeneracy by conditions on small
groups of coordinates.} Fang, Koike, Liu, and Zhao~\cite[Theorem~1.1]{FKLZ} obtain nearly
\(n^{-1/2}\) rates under \skl{sub-exponential tails} and quantitative
two- and three-coordinate covariance conditions, allowing the full
covariance matrix to be singular. The companion paper~\cite{Local}
assumes uniformly positive eigenvalues for principal submatrices of
order at most four. For bounded observations it gives
\(C_aBL^{3/2}/\sqrt n\), and \(C_aBL/\sqrt n\) under aggregate
third-moment cancellation. The present results also cover singular
three- and four-coordinate submatrices. \rev{The two papers therefore
complement each other, and without the local condition the bound loses
only the factor \(\sqrt{D_L}\).}

\skl{Our bounds reach the rate \(n^{-1/2}\), up to logarithmic factors, with
no condition on the covariance matrix, and they lose only the factor
\(\sqrt{D_L}\) against the sharp bounds under nondegeneracy. They come
from an interpolation between the sum and its Gaussian counterpart along
which the derivative of a smoothed rectangle probability is bounded
without reference to the eigenvalues of the covariance matrix.}

\rev{Section~\ref{sec:smoothing} outlines the proof\trim{, derives
the rates from its main estimates, and explains why they reach the
power \(n^{-1/2}\)\lin{, comparing the argument with those of Aronow and
Lopatto and of Gao and Zhang}; it also} introduces the
notation used in the appendices\trim{. Section~\ref{opt:section}
discusses sharpness}, and \skl{Subsection~\ref{sec:lean-certification}}
records the formal verification\skl{ in Lean}.}
Appendices~\trs{\ref{sec:main-proof}}--\ref{app:subexponential}
contain all proofs of the upper bounds. \polish{The Supplement~\cite{Supplement} collects the lower-bound calculations,
exact coordinate reductions, and further investigations.}

\subsection{\texorpdfstring{\sks{Formal verification in Lean}}{Formal verification in Lean}}\label{sec:lean-certification}
\leavevmode\skl{Theorems~\ref{thm:unrestricted}--\ref{opt:consistency} and the two bounds of
Corollary~\ref{cor:subexponential} were certified on 20 and 21~October 2026
in Lean~4.19.0, a proof assistant that checks formal proofs by computer,
with its library mathlib~v4.19.0, and the proofs use only Lean's standard
axioms. The Lean sources, the axiom audits, and the reviews that compare the
formal statements with this manuscript are in the
\href{https://github.com/VC2015/HDCLTs/blob/bbf1ab99c733a398c181ca8c9f5a243a41f258c6/verification/2026-10-21-companion-results/README.md}{verification record}
of the repository, for source revision
\href{https://github.com/VC2015/HDCLTs/commit/6548201a5c5d9c132a53f1d504f7ef70346aa782}{\texttt{6548201a5c5d}}.
The certification refers to these Lean proofs, not to the written
derivations, and the results of the Supplement are not certified.}


% Source: proof_interface
\section{\texorpdfstring{\rs{Outline of the proof}}{Outline of the proof}}
\label{sec:smoothing}

\polish{We first replace each rectangle indicator by a logistic approximation
and then compare its expectations along a path that preserves covariance.
\rev{The smoothing scale balances two errors, since a coarser approximation
costs more at the endpoints but makes the comparison along the path easier.}}
\lin{Subsection~\ref{sec:outline-smoothing} constructs the logistic
approximation at a smoothing scale \(h>0\) and bounds the error of the
replacement (Lemma~\ref{lem:endpoint}), using a bound for the probability
that a sum of independent vectors falls near the boundary of a rectangle
(Lemma~\ref{lem:concentration}). Subsection~\ref{sec:outline-path}
introduces the path, whose parameter \(\lambda\) runs from \(0\) at \(W\)
to \(1\) at~\(Z\), and reduces the theorems to bounds for the derivative
in \(\lambda\) of the smoothed expectations along it.
Proposition~\ref{an:prop:derivative} bounds this derivative \skl{for
\(\lambda\) near \(0\) and for \(\lambda\) away from \(0\)}, and Subsections~\ref{sec:outline-early}
and~\ref{sec:outline-late} explain the two bounds.
Subsection~\ref{sec:rate} integrates them and chooses~\(h\), which proves
the theorems, and Subsection~\ref{sec:why-rate} compares the argument
with those of Aronow and Lopatto and of Gao and Zhang and explains why it
reaches the power \(n^{-1/2}\).

The appendices give the proofs in the same order, from the top down.
Appendix~\ref{sec:main-proof} derives Proposition~\ref{prop:moving-cutoff}
and the theorems from Lemma~\ref{lem:endpoint} and
Proposition~\ref{an:prop:derivative}. Appendix~\ref{appx-sec:smoothing}
proves Lemma~\ref{lem:endpoint} and the properties of the logistic
approximation that the later proofs use, and
Appendix~\ref{appx-sec:support} proves
Proposition~\ref{an:prop:derivative}.
Appendix~\ref{appx-app:concentration} proves
Lemma~\ref{lem:concentration}, and Appendix~\ref{app:subexponential}
treats sub-exponential coordinates.}

\subsection{\texorpdfstring{\rs{The logistic approximation}}{The logistic approximation}}
\label{sec:outline-smoothing}

To construct this approximation, write the rectangle as
\(A=\prod_{j=1}^p[\ell_j,u_j]\), where \(\ell_j\le u_j\),
allowing infinite endpoints.
Each finite endpoint imposes one coordinate bound. If \(e_j\) denotes
the \(j\)th standard basis vector in \(\R^p\), the upper and lower
bounds can both be written with the same inequality sign:
\[
 x_j\le u_j
 \quad\Longleftrightarrow\quad e_j\cdot x\le u_j,
 \qquad
 x_j\ge\ell_j
 \quad\Longleftrightarrow\quad (-e_j)\cdot x\le-\ell_j.
\]
Omit bounds with infinite endpoints and enumerate the remaining
inequalities by \(\alpha=1,\ldots,m\), where \(m\le2p\). We can
then write
\[
 A=\{x\in\R^p:\nu_\alpha\cdot x\le a_\alpha
                 \text{ for }\alpha=1,\ldots,m\}.
\]
Here \(\nu_\alpha\) is the outward normal to the corresponding
coordinate face, and \(a_\alpha\) is its threshold: an upper bound
uses \((\nu_\alpha,a_\alpha)=(e_j,u_j)\), and a lower bound uses
\((\nu_\alpha,a_\alpha)=(-e_j,-\ell_j)\).
Thus the normals are distinct members of
\(\{\pm e_1,\ldots,\pm e_p\}\).

For a smoothing scale \(h>0\), define the logistic approximation to
the indicator function of \(A\):
\begin{equation}
 F_{A,h}(x)=\left(1+\sum_{\alpha=1}^m
          e^{(\nu_\alpha\cdot x-a_\alpha)/h}\right)^{-1},
 \qquad M_h=\sup_A|\E F_{A,h}(W)-\E F_{A,h}(Z)|.
 \label{eq:selector}
\end{equation}
With one inequality, \(F_{A,h}\) is the usual logistic approximation
to a step function. With several inequalities, it has the multinomial
logistic form, with one score for each coordinate bound and a baseline
score of zero. McFadden~\cite[p.~1001]{McFadden1989} uses this form
to smooth indicators of simultaneous inequalities.
\rev{Gao and Zhang~\cite[eq.~(5)]{GaoZhang} use the same multinomial
logistic (softmax) form, with a baseline score, for maxima of high-dimensional sums.}
The quantity \(\nu_\alpha\cdot x-a_\alpha\) is positive exactly
when the corresponding bound is violated. The exponential terms
therefore make \(F_{A,h}\) small outside the rectangle; as \(h\)
decreases to zero, \(F_{A,h}\) approaches its indicator away from
the boundary. The quantity \(M_h\) is the largest difference between
the expectations of these logistic approximations under \(W\) and \(Z\).
Its supremum ranges over all the coordinate bounds just described.

To relate this approximation to the smooth-maximum construction
in~\cite{CCK2013}, observe that, for \(m\ge1\),
\[
 F_{A,h}(x)=g_h\!\left(h\log\sum_{\alpha=1}^m
 e^{(\nu_\alpha\cdot x-a_\alpha)/h}\right),
 \qquad g_h(t)=\frac{1}{1+e^{t/h}}.
\]
The inner function smooths the maximum violation, while \(g_h\)
smooths the indicator of \(({-\infty},0]\). Thus we use a logistic
cutoff with the same scale \(h\) as the smooth maximum; the
construction in~\cite{CCK2013} allows a separate cutoff and a
separate scale at each step.

\lin{An anti-concentration inequality bounds the probability that a
random vector falls in \skl{a boundary layer, that is, within a given
distance of the boundary} of a set. The
following one}, stated
in~\cite[Lemma~2.1]{AL}, controls both the \lin{error of the logistic
approximation and its} derivatives near the boundary. For
\(A=\prod_j[\ell_j,u_j]\), write
\(A^{\rs t}=\prod_j[\ell_j-\rs t,u_j+\rs t]\), with reversed intervals empty.
\begin{lemma}[Rectangle concentration]\label{lem:concentration}
Let \(S=\sum_i\zeta_i\) have independent centered summands with
\(\norm{\zeta_i}_\infty\le\eta\) and
\(\Var(S_j)=v>0\) for every coordinate. For an independent shift
\(R\), a rectangle \(A\), and \(\rs t>0\),
\[
 \Pp(S+R\in A^{\rs t}\setminus A^{-\rs t})
 \le C\min\{1,(\rs t+\eta)\sqrt{\log(ep)/v}\}.
\]
For a centered Gaussian base with equal coordinate variances \(v\),
the same bound holds with \(\eta=0\), including singular covariance.
\end{lemma}
\polish{Appendix~\ref{appx-app:concentration} gives two proofs of the
bounded-sum assertion: a conditioning argument using the Rademacher
anti-concentration inequality in~\cite[Lemma~J.4]{CCKK2022}, and
Aronow--Lopatto's resampling proof~\cite[Lemma~2.1]{AL}, based on
Kane's noise-sensitivity bound~\cite{Kane}\lin{, which controls how
often the indicator of a set changes when a small random fraction of its
inputs is changed}. \rev{\lin{The name of Lemma~\ref{lem:concentration}
refers to L\'evy's concentration function, the largest probability of a
set of given width, which anti-concentration inequalities bound.}} The
Gaussian assertion is
Nazarov's inequality; see~\cite[Lemma~A.1]{CCK2017}. The second proof
also recovers it by a finite-dimensional Gaussian limit.}

\begin{lemma}[Smoothing error]\label{lem:endpoint}
Under~\eqref{eq:setup}, for \(h\ge\delta\),
\begin{equation}
 \E|F_{A,h}(W)-\ind_A(W)|+
 \E|F_{A,h}(Z)-\ind_A(Z)|\le Ch\sqrt L\,\kappa.
 \label{eq:endpoint-error}
\end{equation}
Consequently,
\begin{equation}
 \rh(W,Z)\le M_h+Ch\sqrt L\,\kappa.
 \label{eq:endpoint-min}
\end{equation}
\end{lemma}
\rev{Without a correlation gap the smoothing error is \(ChL^{3/2}\),
because independent unit
exponentials \(\rs{\omega_0},\rs{\omega_\alpha}\) give the exact mixture}
\begin{equation}
 F_{A,h}(x)=\Pp_{\rs\omega}\trs{(x\in A_{\rs\omega}),\qquad
 A_{\rs\omega}=\{x:}\nu_\alpha\cdot x\le
 a_\alpha+h\log(\rs{\omega_\alpha}/\rs{\omega_0})\text{ for every }\alpha\}.
 \label{eq:mixture}
\end{equation}
The expected largest threshold displacement is \(O(hL)\), and
Lemma~\ref{lem:concentration} supplies the factor \(\sqrt L\).
When \(c>0\), a Hessian estimate for
\(h\log(1+\sum_\alpha e^{(\nu_\alpha\cdot x-a_\alpha)/h})\)
gives the alternative error \(Ch\sqrt L/c\). \rev{The parameter \(\kappa\) records the smaller of the two
estimates.} Lemmas~\ref{appx-lem:endpoint-free}
and~\ref{appx-lem:endpoint-gap} prove the two bounds \lin{for the
smoothing error}.

\begin{revblock}
The comparison along the path uses derivatives of \(F_{A,h}\) of orders
three and four, and they are large only near the boundary of the
rectangle. We abbreviate \(F_{A,h}\) to \(F\) when \(A\) and \(h\) are
fixed, and write
\[
 w_\alpha(x)=e^{(\nu_\alpha\cdot x-a_\alpha)/h},\qquad
 \pi_\alpha=\frac{w_\alpha}{1+\sum_\beta w_\beta},\qquad
 \theta=1-F=\sum_\alpha\pi_\alpha,
\]
so that \(\pi_\alpha\) is the multinomial logistic probability of
face~\(\alpha\) and \(\theta\) is the total probability of the faces. By
the chain rule, a derivative of order \(q\) of \(F\) is \(h^{-q}\) times a
sum, over partitions of the \(q\) directions, of products
\(F\pi_{\alpha_1}\cdots\pi_{\alpha_k}\) multiplied by the face projections
\(\nu_\alpha\cdot z\) of the directions~\cite[Lemma~13]{GaoZhang}. For an
ordered coordinate tuple \(I=(j_1,\ldots,j_q)\), write
\(\partial_I=\partial_{j_1}\cdots\partial_{j_q}\). Summing over faces and
tuples gives the \emph{derivative envelope}
\[
 \sum_I\sup_{\norm z_\infty\le h}|\partial_IF(x+z)|
 \le C_qh^{-q}F(x)\theta(x),
\]
which is~\eqref{appx-eq:local-envelope} in
Appendix~\ref{appx-sec:smoothing}. The product \(F\theta\) is at most
\(1/4\), and it is small except within distance of order \(h\) of the
boundary. Moving every threshold outward and inward by \(h\) gives
\(F\theta\le\frac e2(F_{A^h,h}-F_{A^{-h},h})\), and by the mixture
representation the right side is an average of indicators of boundary
\skl{layers} of width \(2h\). Lemma~\ref{lem:concentration} therefore gives
\[
 \E F(U)\theta(U)\le Ch\sqrt L
\]
whenever \(U\) has an independent component to which
Lemma~\ref{lem:concentration} applies with \(v\ge1/2\) and
\(\eta\le h/2\), as Lemma~\ref{appx-lem:envelope} states precisely. On
average, derivatives of order \(q\) are thus of size
\(C_qh^{1-q}\sqrt L\), smaller by the factor \(h\sqrt L\) than the
uniform bound \(h^{-q}\).
\end{revblock}

\subsection{\texorpdfstring{\rs{The interpolation path}}{The interpolation path}}
\label{sec:outline-path}

\polish{For the smooth comparison,} \rev{we use the Bernoulli-thinning
path of Aronow and Lopatto~\cite[\skl{Section~4, the display after eq.~(16)}]{AL}, which removes the bounded
summands at random while a Gaussian vector supplies the lost covariance.}
\polish{Let \(Z\) be independent of the \rs{summands}
and, for \(\lambda\in[0,1]\), set}
\begin{equation}
 \rs\gamma=1+\lambda,\qquad \sks{\varpi}=\frac{1-\lambda}{(1+\lambda)^2},\qquad
 U_\lambda=\rs\gamma\sum_i \sks{\mathbb B}_{i,\lambda}\xi_i+\sqrt\lambda Z,
 \qquad \sks{\mathbb B}_{i,\lambda}\sim\operatorname{Bernoulli}(\sks{\varpi}),
 \label{eq:path}
\end{equation}
with independent Bernoulli variables, \rs{summands}, and Gaussian vector.
Since \(\sks{\varpi}\rs\gamma^2=1-\lambda\), this path has covariance \(\Sigma\),
with \(U_0\stackrel d=W\) and \(U_1\stackrel d=Z\).
\lin{The amplitude \(\gamma=1+\lambda\) is chosen so that the coefficient
\(\sks{\varpi}\gamma^3=1-\lambda^2\) of the third moments of the summands has zero
derivative at the bounded endpoint \(\lambda=0\)~\cite[eq.~(17)]{AL}.}
\rev{Unlike Aronow and Lopatto, we add no Gaussian smoothing at the
bounded endpoint; the logistic approximation takes its place.}
\lin{Put \(P_A(\lambda)=\E F_{A,h}(U_\lambda)\). The largest difference
\(M_h\) of smoothed expectations, defined in~\eqref{eq:selector}, is then
\(\sup_A|P_A(1)-P_A(0)|\), and together with Lemma~\ref{lem:endpoint}
this gives the chain of inequalities
\begin{equation}
 \rh(W,Z)\le M_h+Ch\sqrt L\,\kappa,
 \qquad
 M_h\le\sup_A\int_0^1|P_A'(\lambda)|\dd\lambda,
 \label{eq:roadmap}
\end{equation}
which the proof follows. \skl{We bound} the derivative
\(P_A'(\lambda)\), as Proposition~\ref{an:prop:derivative} below does in
two ranges of~\(\lambda\), and \skl{integrate} the bounds, as
Subsection~\ref{sec:rate} does. For this derivative,} a fourth-order
expansion gives
\begin{equation}
 P_A'(\lambda)=-\frac\lambda3\,
 \E D^3F_{A,h}(U_\lambda):\sum_i\E\xi_i^{\otimes3}
 +R_{4,A,h}(\lambda).
 \label{eq:path-summary}
\end{equation}
\rev{Here the colon denotes contraction of all three indices,
\skl{\(R_{4,A,h}(\lambda)\) collects the fourth-order terms described
below,} and every
cubic tensor is evaluated against the same derivative tensor at
\(U_\lambda\). Under \eqref{eq:zero-third} the cubic terms therefore
cancel before any absolute value is taken, and}
\begin{equation}
 P_A'(\lambda)=R_{4,A,h}(\lambda).
 \label{zs:identity}
\end{equation}
Moving derivatives from \rev{leave-one-out} bases\lin{, from which one
summand is deleted,} back
to \lin{the full path} \(U_\lambda\)
creates fourth-order terms, \rev{which are part of}
\(R_{4,A,h}(\lambda)\).
\begin{revblock}
These terms come from differentiating the law of one summand at a time.
With \(M_i=\E\xi_i\xi_i^\top\) and \(\sks{M_i^{(3)}}=\E\xi_i^{\otimes3}\), Taylor's
formula and the relations \((\sks{\varpi}\gamma^2)'=-1\) and
\((\sks{\varpi}\gamma^3)'=-2\lambda\) give the terms
\(-\frac12M_i:D^2F-\frac\lambda3\sks{M_i^{(3)}}:D^3F\), together with a fourth-order
remainder, at the base \(U_\lambda-\gamma \sks{\mathbb B}_{i,\lambda}\xi_i\) from which
summand~\(i\) is deleted. The Gaussian component contributes
\(\frac12\Sigma:D^2F\) by the heat equation\pass{
\(\frac d{d\lambda}\E f(x+\sqrt\lambda Z)
=\frac12\E\,\Sigma:D^2f(x+\sqrt\lambda Z)\)}, and \(\sum_iM_i=\Sigma\), so
the second-order terms cancel up to the difference of the Hessians at the
full and the leave-one-out bases, which is of fourth order because the
deleted summand is centered and independent of the other terms. Moving the
cubic terms to the full path costs one more fourth-order term. Thus
\(R_{4,A,h}(\lambda)\) is a sum over the summands of three fourth-order
derivatives, each taken at a translate \(U_\lambda+v_i\) with
\(\norm{v_i}_\infty\le4\delta\), in directions given by independent copies
of summand~\(i\); Lemma~\ref{appx-an:lem:identity} gives the exact formula.
\end{revblock}

\begin{proposition}[Two time ranges]\label{an:prop:derivative}
There are sufficiently large absolute constants \(\Ac,d_0\) such that,
for \(h\ge\Ac\delta\), uniformly over the thresholds,
\begin{align}
 |P_A'(\lambda)|&\le C\sqrt L\left(
 \frac{\lambda\delta}{h^2}+\frac{\delta^2}{h^3}\right),
 &&0<\lambda<1,\label{an:early}\\
 |P_A'(\lambda)|&\le CL^{3/2}\left(
 \delta+\frac{\delta^2}{\lambda h}\right),
 &&d_0\delta^2L\le\lambda<1.\label{an:late}
\end{align}
The fourth-order remainder separately satisfies
\begin{align}
 |R_{4,A,h}(\lambda)|&\le C\delta^2\sqrt L/h^3,
 &&0<\lambda<1,\label{eq:remainder-early}\\
 |R_{4,A,h}(\lambda)|&\le C\delta^2L^{3/2}/(\lambda h),
 &&d_0\delta^2L\le\lambda<1.\label{an:remainder}
\end{align}
\end{proposition}
The path is continuous in law at both endpoints, and the logistic approximation is
bounded and continuous. Together with the integrable early derivative
bound, this justifies integration up to \(0\) and \(1\).
\begin{revblock}
\subsection{The early bounds}\label{sec:outline-early}

The early bounds \eqref{an:early} and~\eqref{eq:remainder-early} combine
the derivative envelope with the moment bounds
\[
 \sum_i\E|\xi_{i\ps{j_1}}\xi_{i\ps{j_2}}\xi_{i\ps{j_3}}|\le\delta,\qquad
 \sum_i\E|\xi_{i\ps{j_1}}\ps{\cdots}\xi_{i\ps{j_4}}|\le\delta^2,
\]
which hold for all coordinate indices by the envelope
\(\norm{\xi_i}_\infty\le\delta\), the Cauchy--Schwarz inequality, and the
unit variances. The copies of the summands that supply the directions are
independent of the path, so the coefficients separate from the
derivatives. The shifts \(v_i\) have size at most \(4\delta\le h\), so the
derivative envelope bounds the third and fourth derivatives by
\(Ch^{-3}F\theta(U_\lambda)\) and \(Ch^{-4}F\theta(U_\lambda)\), and the
averaged bound \(\E F\theta(U_\lambda)\le Ch\sqrt L\) applies with the
bounded part of the path as the independent component when
\(\lambda\le1/2\), and with the Gaussian part when \(\lambda\ge1/2\). The
cubic term of~\eqref{eq:path-summary} carries the factor \(\lambda\),
which gives the term \(\lambda\delta/h^2\) in~\eqref{an:early}. The
argument follows the proof of Lemma~4.1 of Aronow and Lopatto~\cite{AL},
where the same derivatives are bounded by Gaussian smoothing, at a cost
that grows as \(\lambda_{\min}(\Sigma)\to0\).

\subsection{The late bounds}\label{sec:outline-late}

Once the path carries Gaussian variance \(\lambda\ge d_0\delta^2L\), the
late bounds \eqref{an:late} and~\eqref{an:remainder} move two derivatives
onto the Gaussian density by Gaussian integration by parts, that is, by
Stein's identity \(\E[Gf(G)]=\E\nabla f(G)\) for a standard Gaussian
vector~\(G\). \trim{Each integration by parts replaces a factor
\(1/h\) by a factor of order \(1/\sqrt\lambda\), so the two together
gain the factor \(h^2\) at the price of a factor \(1/\lambda\)}, and the
four \trim{steps} below keep the resulting bounds free
of any factor that depends on~\(\Sigma\), such as its rank or its
smallest positive eigenvalue.

\emph{Coordinates on the covariance support.} Factor
\(\Sigma=\Gamma\Gamma^\top\), where \(\Gamma\) is a \(p\times r\) matrix of
full column rank \(r=\operatorname{rank}\Sigma\). Each summand lies almost
surely in the range of \(\Gamma\), so \(\xi_i=\Gamma\zeta_i^0\) with
\(\sum_i\E\zeta_i^0(\zeta_i^0)^\top=I_r\), and \(Z\) has the law of
\(\Gamma G\) with \(G\sim N(0,I_r)\). The vectors
\(u_\alpha=\Gamma^\top\nu_\alpha\) have unit length, since
\(\norm{u_\alpha}^2\) is a diagonal entry of \(\Sigma\), and in these
coordinates the rectangle becomes the polyhedron
\(\{x\in\R^r:u_\alpha\cdot x\le a_\alpha\text{ for every }\alpha\}\), an
intersection of halfspaces. We keep the notation \(F\), \(w_\alpha\), and
\(\pi_\alpha\) for the corresponding functions of \(x\in\R^r\), which have
the same form with \(u_\alpha\) in place of~\(\nu_\alpha\). The summands enter
only through the \emph{face norm} \(\norm x_U=\max_\alpha|u_\alpha\cdot x|\),
which satisfies \(\norm{\zeta_i^0}_U\le\delta\), and through the identity
\(\sum_i\E\zeta_i^0(\zeta_i^0)^\top=I_r\). Their Euclidean norms, which are
large when \(\Sigma\) has small positive eigenvalues, are never used.
This reduction follows Gao and Zhang~\cite[Section~2.3]{GaoZhang}.

\emph{Integration by parts against the covariance of the summands.} We write the
path in support coordinates as \(U=V_0+\sigma G\), with \(\sigma^2\) of
order \(\lambda\), where the base \(V_0\) is the bounded part of the path,
together with half of its Gaussian part when \(\lambda>1/2\)\trim{,
since the bounded part alone has variance \(1-\lambda\), which is small
near \(\lambda=1\)}. For a
polyhedron \(Q\) with normals among the \(u_\alpha\), let \(Q^t\) move
every face outward by \(t\). Lemma~\ref{lem:concentration} shows that
\(V_0\) is a \emph{concentrated base}, in the sense that
\(\Pp(V_0\in Q^t\setminus Q^{-t})\le C(t+\delta)\sqrt L\) for every such
\(Q\) and every \(t>0\). \trim{Every Gaussian computation below, such
as an integration by parts, is made with \(V_0\) fixed, so that \(V_0\)
acts as a random translation of the Gaussian vector \(\sigma G\), and
averaging the resulting bounds over~\(V_0\) supplies the small factors.
For a fixed translation, the Gaussian part
alone would bound the probability of a boundary layer of width \(h\)
only by \(Ch\sqrt L/\sigma\), since its scale is \(\sigma\), whereas
Lemma~\ref{lem:concentration} bounds it by \(C(h+\delta)\sqrt L\) for
\(V_0+\sigma G\); the same average gains \lin{a further factor
\(\sigma\) later in this subsection}.} A derivative in a direction \(z\) is
\(\sigma^{-1}\) times a derivative in \(G\), so two integrations by parts
give, for an independent copy \(\zeta_i\) of \(\zeta_i^0\) and directions
\(z_1,\ldots,z_s\),
\[
 \E D^{s+2}F(U)[\zeta_i,\zeta_i,z_1,\ldots,z_s]
 =\sigma^{-2}\E\bigl[\tr\{(GG^\top-I_r)\zeta_i\zeta_i^\top\}\,
 D^sF(U)[z_1,\ldots,z_s]\bigr].
\]
A bound for each summand through the trace of \(GG^\top-I_r\) would cost a
factor \(r\). We sum over the summands first, \trim{using the identity}
\(\sum_i\E\zeta_i\zeta_i^\top=I_r\), and use the duality
\(|\tr(\Xi M)|\le\norm\Xi_\op\norm M_*\) between the operator norm and the
\emph{nuclear norm} \(\norm M_*\), the sum of the absolute values of the
eigenvalues of a symmetric matrix \(M\). Here \(\Xi\) collects the summed
coefficients of the summands, and \(M\) is a Gaussian average that does not
depend on the summand. What is left is the nuclear norm
of Gaussian averages of \(GG^\top-I_r\) against products of the logistic
weights.

\emph{A common base.} The fourth-order terms are evaluated at translates
\(U+v_i\) that depend on the summand, so the weights cannot be taken out
of the sum over the summands directly. With \(N=1+\sum_jw_j\), the identity
\(N(x+v)=N(x)\{1+z_v(x)\}\), where
\(z_v=\sum_j\pi_j(e^{u_j\cdot v/h}-1)\), gives
\begin{equation}
 F\pi_{\alpha_1}\cdots\pi_{\alpha_k}(x+v)
 =F\pi_{\alpha_1}\cdots\pi_{\alpha_k}(x)\,
 e^{\sum_au_{\alpha_a}\cdot v/h}\{1+z_v(x)\}^{-k-1}.
 \label{eq:weight-translation}
\end{equation}
For \(\norm v_U\le4\delta\) we have \(|z_v|\le e^{4\delta/h}-1\), so the
geometric series in \(z_v\) converges absolutely when \(h\ge\Ac\delta\).
This exact expansion moves every factor that depends on the summand into the
matrix \(\Xi\) and leaves products of weights at the common base~\(U\);
Lemma~\ref{appx-an:lem:restoration} carries it out.

\emph{The matrix and surface estimates.} For \(\delta\sqrt L\le\sigma\le1\),
Lemma~\ref{appx-an:lem:matrix} bounds the nuclear norms by
\[
 \E\sum_{\alpha_1,\ldots,\alpha_k}
 \norm{\E_G[(GG^\top-I_r)F\pi_{\alpha_1}\cdots\pi_{\alpha_k}(U)]}_*
 \le C^{k+1}hL^{3/2}.
\]
Each weight \(w=F\pi_{\alpha_1}\cdots\pi_{\alpha_k}(V_0+\sigma\,\cdot)\) is
\emph{log-concave}, that is, has a concave logarithm, because \(\log w\)
is a linear function minus \((k+1)\log N\), and \(\log N\) is convex.
Pr\'ekopa's theorem, which says that Gaussian convolution preserves
log-concavity, then shows that the \emph{tilted} Gaussian law, with
density proportional to \(w\) times the standard Gaussian density, has
covariance \(K\preceq I_r\)\pass{, that is, \(I_r-K\) is positive
semidefinite}. With \(\ps{p_w}=\E w(G)\) and tilted mean \(\mu\), it
follows that
\[
 \norm{\E[(GG^\top-I_r)w(G)]}_*
 =\norm{\ps{p_w}(K-I_r+\mu\mu^\top)}_*
 \le2\ps{p_w}\norm\mu^2-\E[\Hh(G)w(G)],
 \quad \Hh(g)=\norm g^2-r.
\]
\trim{The function \(\Hh\) enters as the trace of \(gg^\top-I_r\), so
that \(\E[\Hh(G)w(G)]=p_w\tr(K-I_r+\mu\mu^\top)\), and the inequality
follows because \(I_r-K\) is positive semidefinite. This trace term may
have either sign. If the weights, summed over the tuples of faces,
formed a partition of unity, as \lin{the indicators of the cells of a
partition of~\(\R^r\)} do, the trace terms would
add up to \(\E\Hh(G)=0\). The logistic weights add up to \(F\theta^k\)
instead, and the trace term is where the geometry of the polyhedron
enters.} An entropy inequality for tilted means, taken relative to one face at a
time (Lemma~\ref{appx-an:lem:mean}), and the averaged bound on
\(F\theta\) control the first term. The second term is a derivative in a
common shift of all thresholds, and it becomes \(h/\sigma\) times the
Gaussian surface integral
\[
 \surf(Q)=\int_{\partial Q}\Hh(g)\varphi_r(g)\dd\mathcal H^{r-1}(g)
\]
of a randomly translated polyhedron, where \(\varphi_r\) is the standard
Gaussian density on \(\R^r\) and \(\mathcal H^{r-1}\) is surface
measure\trim{; moving every face of a polyhedron \(Q\) outward at unit
speed changes the integral of \(\Hh\varphi_r\) over \(Q\) at the rate
\(\surf(Q)\)}.
Gao and Zhang~\cite[Proposition~4]{GaoZhang} prove
\(|\surf(Q)|\le C\{\log(2m)\}^{3/2}\) for every polyhedron with \(m\)
faces, and Lemma~\ref{appx-an:lem:surface} shows that averaging over a
concentrated base gains the factor \(\sigma\), in the sense that
\[
 \E\bigl|\surf\bigl((Q-V_0)/\sigma\bigr)\bigr|
 \le C\bigl(\sigma L^{3/2}+\delta L^2\bigr)\le C\sigma L^{3/2},
 \qquad \sigma\ge\delta\sqrt L.
\]
Its proof uses the Gauss--Green formula\pass{, that is, the divergence
theorem,} within each facet, a face of
dimension \(r-1\), to write \(\surf\) as a sum of facet terms and paired
terms on the ridges, the faces of dimension \(r-2\) where two facets
meet. It then bounds each term by the Gaussian mass of a thin outward
tube\skl{ of width \(1/R\) for a facet or ridge at distance at most \(R\)
from the origin, the sizes \(R\) being summed over dyadic scales in
Subsection~\ref{appx-an:sec:surface}. Such a tube} around the rescaled
polyhedron is a boundary
\skl{layer} of width \(\sigma/R\) for \(V_0+\sigma G\), whose probability
Lemma~\ref{lem:concentration} controls.

In the resulting bound, the integration by parts contributes the factor
\(\sigma^{-2}\), the \(s\) remaining directions contribute
\(\delta^sh^{-s}\) through the derivative formula, and the matrix estimate
contributes \(hL^{3/2}\). For the cubic term, with \(s=1\) and the
coefficient \(\lambda\), this gives \(C\delta L^{3/2}\), and for the
fourth-order terms, with \(s=2\), it gives \(C\delta^2L^{3/2}/(\lambda h)\);
these are \eqref{an:late} and~\eqref{an:remainder}. The condition
\(\lambda\ge d_0\delta^2L\) ensures \(\sigma\ge\delta\sqrt L\).
\trim{Subsection~\ref{sec:why-rate} explains why the factor \(h\) in the
matrix estimate decides the power of the sample size.}
\end{revblock}
\polish{Appendix~\ref{appx-sec:support} gives the exact remainder kernels and proves
these estimates on the covariance support.} \revnote{Proposition~1 is
now the only statement of these bounds; the appendix proves it without
restating it.}


% Source: rate_explanation
\subsection{\texorpdfstring{\trs{Integrating the comparison and choosing the scale}}{Integrating the comparison and choosing the scale}}\label{sec:rate}
\rev{The early and late fourth-order bounds\trim{,
\(C\delta^2\sqrt L/h^3\) and \(C\delta^2L^{3/2}/(\lambda h)\),} of
Proposition~\ref{an:prop:derivative} agree at the time
\(\tau_h=Lh^2\). Integrating the early bounds before this time and the
late bounds after it gives the following \trim{bound for the largest
difference \(M_h=\sup_A|\E F_{A,h}(W)-\E F_{A,h}(Z)|\) of smoothed
expectations, defined in~\eqref{eq:selector}, since
\(M_h\le\sup_A\int_0^1|P_A'(\lambda)|\dd\lambda\)
by~\eqref{eq:roadmap}}. Its proof is in Appendix~\ref{sec:main-proof}.}
\lin{We call a finite system of coordinate inequalities
\(\nu_\alpha\cdot x\le a_\alpha\) a threshold system, and call it
inconsistent if its inequalities have no common solution; such systems
arise in the proofs when thresholds are shifted.}

\begin{revblock}
\begin{proposition}[Smoothed comparison]\label{prop:moving-cutoff}
\label{zs:smooth-proposition}
There are absolute constants \(\mathsf A,C\) such that, under
\eqref{eq:setup}, for every \(h\ge\mathsf A\delta\),
\begin{equation}
 M_h\le CL^{3/2}\left[
 \delta+\frac{\delta^2}{h}\log_+\frac1{Lh^2}\right].
 \tag{S*}\label{eq:smooth-moving}
\end{equation}
If, in addition, the aggregate third-moment tensor vanishes,
as in~\eqref{eq:zero-third}, then
\begin{equation}
 M_h\le\frac{C\delta^2L^{3/2}}h
 \left[1+\log_+\frac1{Lh^2}\right].
 \label{zs:smooth}
\end{equation}
Both estimates are uniform over all finite threshold systems formed from
coordinate inequalities, including inconsistent systems.
\end{proposition}

In~\eqref{eq:smooth-moving}, the cubic term contributes the constant
\(C\delta L^{3/2}\) at every time, while the factor \(1/\lambda\) in the
late fourth-order bound produces the logarithm. Under cancellation the
cubic term is absent, so the early fourth-order contribution, which
\eqref{eq:smooth-moving} absorbs into \(C\delta L^{3/2}\), must be kept
explicitly; it accounts for the additive one in~\eqref{zs:smooth}.
\end{revblock}

\rev{Lemma~\ref{lem:endpoint} bounds the smoothing error by
\(Ch\sqrt L\,\kappa\), and balancing this against the comparison
remainder leads to the scale}
\begin{equation}
 h=\Ac\delta\sqrt{LE_\kappa/\kappa}.
 \label{eq:optimal-scale}
\end{equation}
\trim{Here \(\kappa=\min\{L,c^{-1}\}\), or \(\kappa=L\) when \(c=0\), and
\(E_\kappa=1+\log_+\{\sqrt\kappa/(\delta L)\}\) are the parameters
of~\eqref{eq:parameters}, where \(c\) is the correlation gap.}
Since \(\kappa\le L\) and \(E_\kappa\ge1\), this scale is at
least \(\Ac\delta\). Also
\(1+\log_+(1/(Lh^2))\le CE_\kappa\), so both the smoothing
error and the fourth-order comparison error are at most
\(C\delta L\sqrt{\kappa E_\kappa}\).
This yields Theorems~\ref{thm:optimized} and~\ref{thm:zero-skew};
Theorem~\ref{thm:unrestricted} follows from \(\kappa\le L\) and
\(E_\kappa\le D_L\). Appendix~\ref{sec:main-proof} gives the complete proof, including
\lin{the case \(Lh^2\ge1\), in which the late range is empty,} and the
admissibility of the smoothing scale.

\subsection{\texorpdfstring{\lns{Sources of the argument and the power \(n^{-1/2}\)}}{Sources of the argument and the power n to the -1/2}}
\label{sec:why-rate}\label{sec:sources}
\lin{The interpolation path~\eqref{eq:path} is the Bernoulli-thinning
path of Aronow and Lopatto~\cite[\skl{Section~4, the display after eq.~(16)}]{AL}, and
Lemma~\ref{lem:concentration}, which controls the path near its bounded
endpoint, is their Lemma~2.1; Appendix~\ref{appx-app:concentration}
proves it both by their argument and by a second one. The first two
steps of our exact fourth-order identity
(Lemma~\ref{appx-an:lem:identity}) and the early bounds of
Subsection~\ref{sec:outline-early} follow their proof of Lemma~4.1.
Aronow and Lopatto reach the power \(n^{-1/2}\), but their Gaussian
derivative bounds carry negative powers of \(\lambda_{\min}(\Sigma)\).
The support coordinates of Subsection~\ref{sec:outline-late} follow Gao
and Zhang~\cite[Section~2.3]{GaoZhang}. Along a different path, their
derivative identity~\cite[Lemma~22]{GaoZhang} has the same form as ours,
a cubic term at the full path plus three fourth-order corrections, and
our bound for the corrections (Lemma~\ref{appx-an:lem:restoration})
follows their proof of Lemma~23, which integrates by parts twice in a
Gaussian direction, sums over the summands against their total
covariance, and applies a matrix estimate. Our matrix and surface
estimates, Lemmas~\ref{appx-an:lem:matrix}
and~\ref{appx-an:lem:surface}, adapt their Lemma~12 and Proposition~4,
with the same covariance bound for tilted Gaussian laws and the same
pairing of the boundary terms on each ridge. The logistic
approximation~\eqref{eq:selector} is the multinomial logistic (softmax)
probability of a baseline category with score zero, as used by
\skl{McFadden~\cite[p.~1001]{McFadden1989} and} Gao and
Zhang~\cite[eq.~(5)]{GaoZhang}, and its derivative formula and
\skl{mixture representation~\eqref{eq:mixture}, proved in
Subsection~\ref{appx-sec:mixture}, correspond to} their Lemmas~13 and~14.
Gao and Zhang avoid the
eigenvalues of \(\Sigma\) and obtain the power~\(n^{-1/3}\).}

\trim{The bound of Theorem~\ref{thm:optimized} is the sum of the
smoothing error \(Ch\sqrt L\,\kappa\) of Lemma~\ref{lem:endpoint}, the
cubic contribution \(C\delta L^{3/2}\), which does not depend on~\(h\),
and the integrated fourth-order terms. The power of \(n\) is decided by
the balance between the first and the last. \lin{Gao and Zhang, who
improved the power \(n^{-1/4}\) of~\cite{CCKK2022} to \(n^{-1/3}\),
interpolate in two stages, first from the sum to an intermediate vector
that has a fixed Gaussian component and the same moments through order
three, and then from that vector to the Gaussian
limit~\cite[Section~4]{GaoZhang}. Their bound has the same three
errors~\cite[eq.~(20)]{GaoZhang}. In our notation, with \(h\) in place
of their \(1/\beta\) and \(\delta\) in place of their envelope
\(b_{d,n}\), the fourth-order terms of their first stage are bounded,
when that stage carries Gaussian variance \(\lambda/2\), in the proof of
their Proposition~2, by a constant times}
\[
 \frac{\delta^2L}{\lambda h^2}.
\]
\lin{Integrated from a small time \(\tau\) to one}, this gives
\(\delta^2L/h^2\) times \lns{\(\log(1/\tau)\)}, which balances the
smoothing error \(h\sqrt L\) when \(h\) is of order \(\delta^{2/3}\), up
to logarithmic factors. This is their choice of~\(\beta\)
in~\cite[eq.~(21)]{GaoZhang}, and with \(\delta=B/\sqrt n\) it gives the
power \(n^{-1/3}\). Our late bound~\eqref{an:remainder} is smaller by
the factor \(h\sqrt L\). Its integral, \(\delta^2L^{3/2}/h\) times a
logarithm, balances the smoothing error when \(h\) is of order
\(\delta\), up to logarithmic factors, as in~\eqref{eq:optimal-scale},
and this gives the power \(n^{-1/2}\).

\skl{The derivatives of \(F_{A,h}\) are large only within distance of order \(h\)
of the boundary of the rectangle, and the factor \(h\sqrt L\) is the
probability of that boundary layer, which the} anti-concentration inequality
of Lemma~\ref{lem:concentration} bounds \skl{by}
\(Ch\sqrt L\); this is the bound \(\E F\theta(U)\le Ch\sqrt L\) of
Subsection~\ref{sec:outline-smoothing}. In the early range the derivative
envelope gives the factor directly (Subsection~\ref{sec:outline-early}).
In the late range two derivatives are moved onto the Gaussian density,
and the factor must survive this step. Gao and Zhang write the products
of logistic weights that remain after the integrations by parts as
probabilities of convex cells that partition the
space~\cite[Lemmas~13 and~14]{GaoZhang}, and their bound for the sum over
the cells~\cite[Lemma~12]{GaoZhang} is of order \(L\), uniformly over
partitions. A partition has total probability one, so this bound\skl{,
being uniform over partitions, does not carry the factor \(h\sqrt L\), and
we see no way to recover the factor from it}.

\skl{The factor survives in our argument because the matrix estimate is
applied to the logistic weights themselves and because the trace term,
which does not cancel, is averaged over the concentrated base. The total
mass of the weights satisfies} \(\E F\theta^k(U)\le Ch\sqrt L\), and the term of \skl{the
proof of the matrix estimate} that involves
the tilted means inherits this factor directly. \skl{The} trace term
\(\E[\Hh(G)w(G)]\) does not cancel, because the weights do not form a
partition, and it is bounded by \(h/\sigma\) times an average of Gaussian
surface integrals of translated polyhedra. The uniform bound
\(C\{\log(2m)\}^{3/2}\) of Gao and Zhang~\cite[Proposition~4]{GaoZhang}
for these integrals would leave the factor \(h/\sigma\), and the late
bound would become \(C\delta^2L^{3/2}/(\lambda^{3/2}h)\). The integral of
this bound \skl{over the late range} is not logarithmic\skl{. At the time
\(\tau_h=Lh^2\) of Subsection~\ref{sec:rate} the balance with the
smoothing error gives the power \(n^{-1/3}\) of Gao and Zhang, and only a
time re-optimized to the order \(L^{2/3}h^{4/3}\) gives \(n^{-3/8}\), up to
logarithmic factors.} The averaged surface estimate gains the factor \(\sigma\) by averaging over
the concentrated base, which restores the factor \(1/\lambda\) and with
it the logarithm. Gao and Zhang use their uniform bound for the cubic
term \lin{of their second stage}~\cite[eq.~(19)]{GaoZhang}, where it suffices,
because that term involves no power of~\(h\).

\lin{The translation formula~\eqref{eq:weight-translation} and the
exact identity of Lemma~\ref{appx-an:lem:identity} replace
approximations in the argument of Gao and Zhang. The first brings every
derivative from the translated bases \(U+v_i\) back to the common base
\(U\) as an absolutely convergent series with no remainder term, in place
of their independent copies and high-order Taylor
expansion~\cite[Lemmas~19 and~23]{GaoZhang}. The second places the cubic
term at the full path, so that cancellation of third moments removes it
before any estimate is taken, which gives
Theorem~\ref{thm:zero-skew}.}}


% Source: sharpness_and_discussion
\section{Sharpness and remaining questions}\label{opt:section}
The dependence on sample size cannot improve beyond \(n^{-1/2}\)
under the present assumptions. A normalized Rademacher sum has an atom
of order \(n^{-1/2}\), while the Gaussian probability of an interval
shrinking to this atom tends to zero. This obstruction already occurs
in one dimension and survives zero skewness\lin{, that is, the
cancellation~\eqref{eq:zero-third} of third moments}.

Dimension introduces another obstruction, already visible when the
coordinates are independent. \pass{Chernozhukov, Chetverikov, and
Koike~\cite[Proposition~2.1]{CCK} showed with i.i.d.\ two-point
coordinates that the factor \((\log p)^{3/2}\) in the rate cannot be
improved, under further growth conditions on \(B_n\) and \(p_n\) that the
following theorem does not need.}
\begin{theorem}[Lower scale]\label{opt:skew-lower}
Suppose \(p_n\to\infty\), \(B_n\ge2\), and
\(r_n=B_n(\log p_n)^{3/2}/\sqrt n\to0\). There are i.i.d.
centered observations with independent coordinates, covariance
\(I_{p_n}\), and envelope \(B_n\), for which
\begin{equation}
 \liminf_n\frac{\rh(W_n,Z_n)}{r_n}\ge\frac{\sqrt2}{4e}.
 \label{opt:lower-constant}
\end{equation}
Here \(W_n=n^{-1/2}\sum_iX_{n,i}\),
\(Z_n\sim N(0,I_{p_n})\), and the envelope is attained.
\end{theorem}
The example uses the centered, unit-variance two-point law
\[
 \Pp(X=B_n)=\frac1{B_n^2+1},\qquad
 \Pp(X=-B_n^{-1})=\frac{B_n^2}{B_n^2+1}.
\]
At a normal maximum threshold of order \(\sqrt{\log p_n}\), its
skewness changes each relative tail by order \(r_n\). Taking the
product of the coordinate probabilities preserves this discrepancy.
\polish{The Supplement~\cite[Section~S.4.1, Theorem~S.3 and Lemma~S.2]{Supplement}
proves a moderate-tail estimate uniform in the envelope and includes the
lattice corrections. This extends the regime of the two-point lower
construction in~\cite[Proposition~2.1]{CCK}, with the law parameterized here
by its attained envelope.}

\rev{These lower examples have correlation gap one, so they belong to
every class with a fixed positive gap.} To express uniformity,
let \(M_n(c_0;B,p)\) be the supremum of rectangle error over the
independent centered observations with envelope at most \(B\), unit
average marginal variances, and correlation gap at least \(c_0\).
Write \(M_n^0(c_0;B,p)\) for the same supremum restricted to
\pass{observations whose aggregate third-moment tensor vanishes}.
For fixed \(c_0\in(0,1]\), Theorem~\ref{thm:optimized} and the lower
example give
\begin{equation}
 M_n(c_0;B_n,p_n)\asymp_{c_0,\rs\Lambda}
 \frac{B_n(\log p_n)^{3/2}}{\sqrt n}
 \label{opt:matching-minimax}
\end{equation}
whenever \(p_n\to\infty\), \(B_n\ge2\), the displayed rate tends
to zero, and \(1+\log_+(\sqrt n/B_n)\le\rs\Lambda\log(ep_n)\).
\polish{This includes polynomial dimension and growing envelopes subject to
the vanishing-rate condition. The Supplement~\cite[Section~S.4.2,
Theorem~S.4]{Supplement} gives the general consistency argument below.}

\begin{theorem}[Uniform consistency at a fixed gap]
\label{opt:consistency}
For fixed \(c_0\in(0,1]\), arbitrary dimension sequences \(p_n\ge1\),
and envelopes \(B_n\ge2\),
\begin{equation}
 M_n(c_0;B_n,p_n)\longrightarrow0
 \quad\Longleftrightarrow\quad
 \frac{B_n\{\log(ep_n)\}^{3/2}}{\sqrt n}\longrightarrow0.
 \label{opt:consistency-iff}
\end{equation}
For each fixed \(B\ge1\),
\[
 M_n^0(c_0;B,p_n)\longrightarrow0
 \quad\Longleftrightarrow\quad \log p_n=o(\sqrt n).
\]
\end{theorem}

For the zero-skewness assertion, independent Rademacher coordinates
locate the critical dimension at
\(\log p\asymp\sqrt n\). If \(\log p_n\sim a\sqrt n\), their
maximum has a limiting discrepancy
\(e^{-\exp(-a^2/3)}-e^{-1}>0\) at a threshold chosen so that the
Gaussian upper-tail count tends to one. \polish{The fourth-order tail calculation and the necessity argument for
arbitrary dimension sequences are in the Supplement~\cite[Section~S.3,
Proposition~S.2 and the following proof]{Supplement}.} \rev{The sharpness of the
threshold says nothing about the error rate below it. In the range
\(n^{1/3}\ll L\ll\sqrt n\), known lower bounds have scale
\(L^2/n\), while the fixed-gap zero-skewness upper bound is larger.}
The lattice obstruction must also be retained in any improved
unsmoothed error bound.

\rev{For unrestricted covariance the upper bound is optimal up to
logarithmic factors, and where the lower scale \(\delta L^{3/2}\) is
attained, the upper and lower bounds differ by at most the factor
\(\sqrt{D_L}\). We do not know whether this factor is necessary, nor how
the optimal bound depends on a vanishing correlation gap. The companion
note~\cite{Obstruction} shows that natural pointwise strengthenings of
the fourth-order remainder bound, for a two-scale version of the
smoother, are false, so the factor cannot be removed by recovering a
cancellation lost in a pointwise derivative estimate. It also states an
integrated estimate of the remainder that would suffice, and a second
note~\cite{IntegratedRemainder} proves this estimate when the covariance
has rank one.}
\revnote{Removed: the claim that a companion note removes the factor at
covariance rank at most two. That note is not available.}

\polish{Exactly repeated or opposite coordinates can be removed without
changing rectangle distance; see the Supplement~\cite[Section~S.5,
Proposition~S.6]{Supplement}. It also gives the rank-one minimax rate
\(C\min\{1,B/\sqrt n\}\) \cite[Section~S.4.3, Proposition~S.4]{Supplement}.
\rev{The proof for arbitrary rank does not use these reductions.}}

\rev{Corollary~\ref{cor:subexponential} complements the companion
paper~\cite{Local}, whose sub-exponential bound is sharper in logarithms
when all principal covariance submatrices of order at most four are
uniformly nondegenerate. Without that restriction,
Corollary~\ref{cor:subexponential} still gives the power \(n^{-1/2}\) up
to logarithms. Appendix~\ref{app:subexponential} compares the
assumptions and rates in detail.}


% Source: references
\begin{thebibliography}{99}
\phantomsection
\addcontentsline{toc}{section}{References}
\bibitem{AL}
P.~M. Aronow and P. Lopatto (2026).
\emph{Optimal central limit theorem for bounded random variables in high dimensions}.
Manuscript, arXiv:2609.07067v1.
\url{https://arxiv.org/abs/2609.07067v1}.

\bibitem{CCK2013}
V. Chernozhukov, D. Chetverikov, and K. Kato (2013).
\emph{Gaussian approximations and multiplier bootstrap for maxima of
sums of high-dimensional random vectors}.
The Annals of Statistics \textbf{41}(6), 2786--2819.
\href{https://doi.org/10.1214/13-AOS1161}{doi:10.1214/13-AOS1161}.

\bibitem{CCK2017}
V. Chernozhukov, D. Chetverikov, and K. Kato (2017).
\emph{Central limit theorems and bootstrap in high dimensions}.
The Annals of Probability \textbf{45}(4), 2309--2352.
\href{https://doi.org/10.1214/16-AOP1113}{doi:10.1214/16-AOP1113}.

\bibitem{CCKK2022}
V. Chernozhukov, D. Chetverikov, K. Kato, and Y. Koike (2022).
\emph{Improved central limit theorem and bootstrap approximations in
high dimensions}. The Annals of Statistics \textbf{50}(5), 2562--2586.
\href{https://doi.org/10.1214/22-AOS2193}{doi:10.1214/22-AOS2193}.

\bibitem{Local}
V. Chernozhukov, D. Chetverikov, K. Kato, and Y. Koike,
with AI collaborators (2026).
\emph{High-dimensional Gaussian approximation over rectangles under
local nondegeneracy}. Companion working manuscript\skl{, 21 October 2026}.

\bibitem{Supplement}
\polish{V.~Chernozhukov, D.~Chetverikov, K.~Kato, and Y.~Koike,
with AI collaborators (2026).
\emph{Supplementary material: High-dimensional Gaussian approximations
over rectangles without covariance restrictions}.
Working manuscript, 20 October 2026, revision \texttt{bbf1ab99c733}.
\href{https://github.com/VC2015/HDCLTs/blob/bbf1ab99c733a398c181ca8c9f5a243a41f258c6/research-rectangle-clt-singular/Rectangle_CLT_Supplement.pdf}{Supplement}.}

\bibitem{CCK}
V. Chernozhukov, D. Chetverikov, and Y. Koike (2023).
\emph{Nearly optimal central limit theorem and bootstrap approximations
in high dimensions}.
The Annals of Applied Probability \textbf{33}(3), 2374--2425.
\href{https://doi.org/10.1214/22-AAP1870}{doi:10.1214/22-AAP1870}.
Working version: arXiv:2012.09513v2, May 12, 2021.
\url{https://arxiv.org/abs/2012.09513v2}.

\bibitem{FKLZ}
X. Fang, Y. Koike, S.-H. Liu, and Y.-K. Zhao (2026).
\emph{High-dimensional central limit theorems by Stein's method in
the degenerate case}. The Annals of Applied Probability, to appear.
\skl{arXiv:2305.17365.}
\href{https://research.cuhk.edu.hk/en/publications/high-dimensional-central-limit-theorems-by-steins-method-in-the-d/}{Accepted manuscript (2026)}.

\bibitem{GaoZhang}
Z. Gao and W. Zhang (2026).
\emph{Cubic-root Gaussian approximation under unrestricted covariance}.
arXiv:2608.30221v1, 31 August 2026.
\url{https://arxiv.org/abs/2608.30221v1}.

\bibitem{Kane}
\polish{D.~M. Kane (2014).
\emph{The average sensitivity of an intersection of half spaces}.
Research in the Mathematical Sciences \textbf{1}, Article~13.
\href{https://doi.org/10.1186/s40687-014-0013-6}{doi:10.1186/s40687-014-0013-6}.
Working version: arXiv:1309.2987v3.
\url{https://arxiv.org/abs/1309.2987v3}.}

\bibitem{McFadden1989}
D. McFadden (1989).
\emph{A method of simulated moments for estimation of discrete
response models without numerical integration}.
Econometrica \textbf{57}(5), 995--1026.
\href{https://doi.org/10.2307/1913621}{doi:10.2307/1913621}.

\bibitem{OST}
R. O'Donnell, R.~A. Servedio, and L.-Y. Tan (2019).
Fooling polytopes.
In \emph{Proceedings of the 51st Annual ACM SIGACT Symposium on
Theory of Computing}, 614--625.
\href{https://doi.org/10.1145/3313276.3316321}{doi:10.1145/3313276.3316321}.

\bibitem{Prekopa1973}
\polish{A.~Pr\'ekopa (1973).
\emph{On logarithmic concave measures and functions}.
Acta Scientiarum Mathematicarum (Szeged) \textbf{34}, 335--343.
\url{https://rutcor.rutgers.edu/Prekopa/pdf/SCIENT2.pdf}.}

\bibitem{IntegratedRemainder}
\rev{\emph{Toward an integrated curvature estimate: a proved rank-one case
and an entropy identity for the general problem}.
Companion research note, 24 October 2026.
\href{https://github.com/VC2015/HDCLTs/blob/main/research-rectangle-clt-singular/Rectangle_CLT_Integrated_Remainder.pdf}{\nolinkurl{Rectangle_CLT_Integrated_Remainder.pdf}}.}

\bibitem{Obstruction}
\rev{\emph{A curvature obstruction for two-scale smoothing}.
Companion research note, 23 October 2026.
\href{https://github.com/VC2015/HDCLTs/blob/main/research-rectangle-clt-singular/Rectangle_CLT_Curvature_Obstruction.pdf}{\nolinkurl{Rectangle_CLT_Curvature_Obstruction.pdf}}.}

\end{thebibliography}



\clearpage
\appendix
\numberwithin{equation}{section}
\numberwithin{theorem}{section}
\numberwithin{lemma}{section}
\numberwithin{proposition}{section}
\numberwithin{corollary}{section}
\numberwithin{remark}{section}
\begin{revblock}
\section*{Notation in the appendices}
\addcontentsline{toc}{section}{Notation in the appendices}
The following symbols, most of which Section~\ref{sec:smoothing}
introduces, are used throughout
Appendices~\trs{\ref{sec:main-proof}, \ref{appx-sec:smoothing},}
and~\trs{\ref{appx-sec:support}}.
Appendices~\trs{\ref{appx-app:concentration}} and~E use a few local
symbols, introduced where they occur.

\medskip
\begin{center}
\small
\renewcommand{\arraystretch}{1.15}
\begin{tabular}{@{}p{0.27\textwidth}p{0.68\textwidth}@{}}
\trim{\(\delta,\ L,\ c,\ \kappa,\ E_\kappa\)} & \trim{summand envelope,
\(\log(ep)\), correlation gap, and the parameters
of~\eqref{eq:parameters}}\\
\trim{\(F_{A,h},\ M_h\)} & \trim{logistic approximation of \(A\) at scale
\(h\), and \(\sup_A|\E F_{A,h}(W)-\E F_{A,h}(Z)|\),
see~\eqref{eq:selector}}\\
\trim{\(P_A(\lambda),\ \tau_h\)} & \trim{\(\E F_{A,h}(U_\lambda)\) along
the path~\eqref{eq:path}, and the time \(Lh^2\) at which the early and
late bounds meet}\\
\(\nu_\alpha,\ a_\alpha,\ m\) & \trim{outward normal \(\pm e_j\)} and threshold of
face \(\alpha\); there are \(m\le2p\) finite faces\\
\(A^t\) & the rectangle \(A\) with every finite endpoint moved outward
by \(t\), or inward if \(t<0\)\\
\(F,\ \Psi,\ w_\alpha,\ T\trs{,\ N}\) & logistic approximation, its potential
\(h\log(1+T)\), the face weights, \trim{their} sum
\(T=\sum_\alpha w_\alpha\)\trim{, and the denominator \(N=1+T\) of \(F\)}\\
\(\pi_\alpha,\ \theta,\ \bar\nu,\ \chi\) & normalized weights, total
mass \(1-F\), weighted normal \(\nabla\Psi\), and curvature
\(h\tr D^2\Psi\)\\
\(\omega_0,\ \omega_\alpha\) & independent unit exponentials in the
mixture representation of \(F\)\\
\(\Gamma,\ r,\ u_\alpha,\ \norm\cdot_U\) & support factor
\(\Sigma=\Gamma\Gamma^\top\), its rank, face normals
\(\Gamma^\top\nu_\alpha\), and the face norm
\(\max_\alpha|u_\alpha\cdot x|\)\\
\(\zeta_i^0;\ \zeta_i,\ \widetilde\zeta_i\) & summands in support coordinates; two
independent copies of them\\
\(\gamma,\ \sks{\varpi},\ \sks{\mathbb B}_{i,\lambda}\) & amplitude \(1+\lambda\), selection
probability \((1-\lambda)/(1+\lambda)^2\), and the selection variables\\
\(Y_i,\ V_\lambda,\ U_\lambda,\ H_i\) & selected summands
\(\gamma \sks{\mathbb B}_{i,\lambda}\zeta_i^0\), their sum, the path, and the path with
summand \(i\) deleted\\
\(V_0,\ \sigma,\ G,\ U\) & concentrated base, Gaussian scale, standard
Gaussian vector on \(\R^r\), and \(U=V_0+\sigma G\)\skl{; the coordinate
standard deviations \(\sigma_j\) of Corollary~\ref{cor:subexponential}
carry a subscript}\\
\(\Phi_r,\ \varphi_r,\ H_2,\ \surf\) & standard Gaussian measure and
density on \(\R^r\), \(H_2(g)=\norm g^2-r\), and the Gaussian surface
integral\\
\(\mathsf A,\ \mathsf A_b,\ d_0\) & large absolute constants in
\(h\ge\mathsf A\delta\), \(h\ge\mathsf A_b\delta\), and
\(\lambda\ge d_0\delta^2L\)\\
\(C,\ C_q;\ c_1,\ c_2,\ c_*\) & absolute constants, large and small;
they may change between occurrences, and a subscript records any
dependence. The letter \(c\) alone is the correlation gap.\\
\end{tabular}
\end{center}
\end{revblock}

% Source: complete_integration_proofs
\section{\texorpdfstring{\trs{Proofs of Proposition~\ref{prop:moving-cutoff} and the theorems}}{Proofs of Proposition 2 and the theorems}}\label{sec:main-proof}
\trim{This appendix proves Proposition~\ref{prop:moving-cutoff} and
Theorems~\ref{thm:unrestricted}--\ref{thm:zero-skew} along the lines of
Subsection~\ref{sec:rate}. It uses only Lemma~\ref{lem:endpoint} and
Proposition~\ref{an:prop:derivative}, which
Appendices~\ref{appx-sec:smoothing} and~\ref{appx-sec:support} prove.
Recall that \(P_A(\lambda)=\E F_{A,h}(U_\lambda)\) along the
path~\eqref{eq:path}, so that \(P_A(0)=\E F_{A,h}(W)\) and
\(P_A(1)=\E F_{A,h}(Z)\), and that
\(M_h=\sup_A|P_A(1)-P_A(0)|\) by~\eqref{eq:selector}.}

\begin{proof}[\rev{Proof of~\eqref{eq:smooth-moving} in
Proposition~\ref{prop:moving-cutoff}}]
Increase \(\Ac\) so that \eqref{an:early} and \eqref{an:late} apply and
\(\Ac^2\ge d_0\). Put \(\tau_h=Lh^2\), the time at which their
fourth-order terms agree. Since \(h\ge\Ac\delta\), we have
\(\tau_h\ge d_0\delta^2L\). If \(\tau_h<1\), the later estimate
is therefore available throughout \([\tau_h,1]\). Integrating the
early estimate gives
\begin{align*}
 |P_A(\tau_h)-P_A(0)|
 &\le C\sqrt L\left(
 \frac{\delta\tau_h^2}{h^2}
 +\frac{\delta^2\tau_h}{h^3}\right)\\
 &=C\delta L^{3/2}
 \left(Lh^2+\frac\delta h\right)
 \le C\delta L^{3/2}.
\end{align*}
The later estimate gives
\[
 |P_A(1)-P_A(\tau_h)|
 \le CL^{3/2}\left[
 \delta+\frac{\delta^2}{h}\log\frac1{\tau_h}\right].
\]
The constant cubic contribution integrates to \(C\delta L^{3/2}\);
the fourth-order factor \(1/\lambda\) gives the logarithm. Combining
the two time ranges and taking the supremum over thresholds proves
\eqref{eq:smooth-moving} when \(\tau_h<1\).

If \(Lh^2\ge1\), integrate the early estimate over the whole path:
\[
 |P_A(1)-P_A(0)|
 \le C\sqrt L\left(\frac\delta{h^2}
 +\frac{\delta^2}{h^3}\right)
 =C\delta L^{3/2}\frac{1+\delta/h}{Lh^2}
 \le C\delta L^{3/2}.
\]
The logarithmic term in \eqref{eq:smooth-moving} is zero in this case.
Continuity at the interpolation endpoints permits integration up to
\(0\) and \(1\). All bounds are uniform in the thresholds.
\end{proof}

\begin{proof}[Proof of Theorem~\ref{thm:optimized}]
Combining Proposition~\ref{prop:moving-cutoff} with
\eqref{eq:endpoint-min} yields, for every \(h\ge\Ac\delta\),
\begin{equation}
 \rh(W,Z)\le C\left[
 \delta L^{3/2}
 +\frac{\delta^2}{h}L^{3/2}\log_+\frac1{Lh^2}
 +h\sqrt L\,\kappa\right].
 \label{eq:variational}
\end{equation}
The first term is independent of \(h\). To balance the fourth-order
comparison error against the endpoint error, choose \trim{\(h\) as
in~\eqref{eq:optimal-scale}}. Because \(\kappa\le L\) and \(E_\kappa\ge1\), this scale is at least
\(\Ac\delta\). Moreover,
\[
 \log_+\frac1{Lh^2}
 =2\log_+\frac{\sqrt\kappa}{\Ac\delta L\sqrt{E_\kappa}}
 \le2\log_+\frac{\sqrt\kappa}{\delta L}
 \le2E_\kappa.
\]
The endpoint term in \eqref{eq:variational} equals
\(\Ac\delta L\sqrt{\kappa E_\kappa}\), while its logarithmic term is at
most \((2/\Ac)\delta L\sqrt{\kappa E_\kappa}\).
This proves \eqref{eq:optimized} after taking the minimum with one.
The version with \(D_L\) follows from \(E_\kappa\le D_L\).
\end{proof}


\begin{proof}[Proof of Theorem~\ref{thm:unrestricted}]
Since \(\kappa\le L\), \(E_\kappa\le D_L\), and \(D_L\ge1\),
\[
 L^{3/2}+L\sqrt{\kappa E_\kappa}
 \le 2L^{3/2}\sqrt{D_L}.
\]
Apply Theorem~\ref{thm:optimized}.
\end{proof}


% Source: zero_skewness_proof
\subsection{\texorpdfstring{\sks{Third-moment cancellation: the second part of Proposition~\ref{prop:moving-cutoff} and Theorem~\ref{thm:zero-skew}}}{Third-moment cancellation: the second part of Proposition 2 and Theorem 3}}\label{zs:section}
\trim{This subsection proves the second part of
Proposition~\ref{prop:moving-cutoff} and Theorem~\ref{thm:zero-skew}.}

\begin{proof}[\rev{Proof of~\eqref{zs:smooth} in
Proposition~\ref{prop:moving-cutoff}}]
Use \eqref{zs:identity} to cancel the common-base cubic term before
estimating the derivative. The separate fourth-order bounds
\eqref{eq:remainder-early} and \eqref{an:remainder} then give
\begin{align}
 |P_A'(\lambda)|
 &\le \frac{C\delta^2\sqrt L}{h^3},
 &&0<\lambda<1,\label{zs:early}\\
 |P_A'(\lambda)|
 &\le \frac{C\delta^2L^{3/2}}{\lambda h},
 &&d_0\delta^2L\le\lambda<1.\label{zs:late}
\end{align}
Increase \(\Ac\) so that \(\Ac^2\ge d_0\), and put \(\tau_h=Lh^2\).
If \(\tau_h<1\), then \(\tau_h\ge d_0\delta^2L\). Integrating
\eqref{zs:early} up to \(\tau_h\), and \eqref{zs:late} after it,
yields
\begin{align*}
 |P_A(1)-P_A(0)|
 &\le \frac{C\delta^2\sqrt L}{h^3}\tau_h
       +\frac{C\delta^2L^{3/2}}h\log\frac1{\tau_h}\\
 &=\frac{C\delta^2L^{3/2}}h
       \left[1+\log\frac1{Lh^2}\right].
\end{align*}
If \(Lh^2\ge1\), integrate \eqref{zs:early} on the whole path and
use \(\sqrt L/h^3\le L^{3/2}/h\). Endpoint continuity, already
established after Proposition~\ref{an:prop:derivative}, justifies
integration to \(0\) and \(1\). Taking the supremum over thresholds
proves \eqref{zs:smooth}.
\end{proof}

\begin{proof}[Proof of Theorem~\ref{thm:zero-skew}]
Combine \rev{the second part of} Proposition~\rev{\ref{prop:moving-cutoff}} with the endpoint
bound \eqref{eq:endpoint-min} to obtain, for every \(h\ge\Ac\delta\),
\begin{equation}
 \rh(W,Z)\le C\left\{
 \frac{\delta^2L^{3/2}}h
 \left[1+\log_+\frac1{Lh^2}\right]+h\sqrt L\,\kappa\right\}.
 \label{zs:variational}
\end{equation}
To balance the comparison error against the endpoint error, choose
\trim{\(h\) by}~\eqref{eq:optimal-scale} again\trim{; it} is admissible because \(\kappa\le L\) and \(E_\kappa\ge1\).
As before,
\[
 \log_+\frac1{Lh^2}
 =2\log_+\frac{\sqrt\kappa}{\Ac\delta L\sqrt{E_\kappa}}
 \le2E_\kappa.
\]
The endpoint term in \eqref{zs:variational} equals
\(\Ac\delta L\sqrt{\kappa E_\kappa}\); its comparison term is at most
\[
 \frac{\delta L\sqrt\kappa}{\Ac\sqrt{E_\kappa}}(1+2E_\kappa)
 \le\frac3{\Ac}\delta L\sqrt{\kappa E_\kappa}.
\]
Taking the minimum with one proves \eqref{eq:zero-skew}.

\end{proof}


% Source: smoothing_proofs
\section{\texorpdfstring{\trs{The logistic approximation: proof of Lemma~\ref{lem:endpoint}}}{The logistic approximation: proof of Lemma 2}}\label{appx-sec:smoothing}

\trim{This appendix develops the properties of the logistic
approximation \(F_{A,h}\) of Subsection~\ref{sec:outline-smoothing} that
the proofs use. We first record how Lemma~\ref{lem:concentration}
applies to systems of halfspaces, then derive the mixture and derivative
formulas for \(F_{A,h}\), and finally bound the error that it introduces
at \(W\) and \(Z\)\trim{, which proves Lemma~\ref{lem:endpoint}. The
derivative formulas and the averaged envelope of
Lemma~\ref{appx-lem:envelope} are used again in
Appendix~\ref{appx-sec:support}}.}

\subsection{\texorpdfstring{\trs{Anti-concentration for systems of halfspaces}}{Anti-concentration for systems of halfspaces}}

\trim{For} a fixed system of halfspaces\trim{, the notation
\(A^{\rs t}\) of Subsection~\ref{sec:outline-smoothing}} increases all its
thresholds by \(\rs t\). We also apply \trim{Lemma~\ref{lem:concentration},
which Appendix~\ref{appx-app:concentration} proves,} to the
\trim{vector of face projections}
\((\nu_\alpha\cdot x)_{\alpha\le m}\). Its coordinates retain
the same variance and summand bound. The intersection
\(\{\nu_\alpha\cdot x\le a_\alpha\}_{\alpha\le m}\) is the inverse
image of the lower orthant \(\prod_{\alpha\le m}(-\infty,a_\alpha]\)
in \(\R^m\); comparing thresholds increased and decreased by \(\rs t\)
gives the \skl{boundary layer} in the concentration estimate. The orthant \pass{is} nonempty even when
its inverse image is empty, so \trim{systems of inequalities with no
common solution} are covered. Since \(m\le2p\), \pass{we have} \(\log(em)\le CL\).

\subsection{\texorpdfstring{\trs{Mixture and derivative formulas}}{Mixture and derivative formulas}}\label{appx-sec:mixture}

\trim{We keep the notation of Subsection~\ref{sec:outline-smoothing} for
the faces \(\alpha=1,\ldots,m\) of \(A\), with \trim{outward} normals
\(\nu_\alpha\) and thresholds \(a_\alpha\), for the face weights
\(w_\alpha\), the logistic approximation \(F=F_{A,h}\), the face
probabilities \(\pi_\alpha\), and \(\theta=1-F\), and we put in
addition}
\begin{gather*}
 \trs{T=\sum_\alpha w_\alpha,\qquad \pi_0=F=\frac1{1+T},\qquad
 \Psi=\Psi_{A,h}=h\log(1+T),}\\
 \rs{\bar\nu}=\sum_\alpha\pi_\alpha\nu_\alpha,
 \qquad \chi=h\,\tr D^2\Psi.
\end{gather*}
\trim{Here \(\pi_0\) is the probability of the baseline category, the
vector} \(\rs{\bar\nu}=\nabla\Psi\) is \trim{the} weighted normal\trim{, and}
the scalar \(\chi\) measures the curvature used in the gap-sensitive
endpoint estimate below. We use \trim{these formulas} for arbitrary
coordinate thresholds, including those with an empty intersection,
because the shifts below need not preserve the ordering of the
endpoints. In such cases, \(F_{A,h}\) refers to the specified
inequalities, and its denominator is positive. If \(m=0\), the logistic
approximation is identically one and all comparisons are trivial; the
analytic arguments below take \(m\ge1\).

\trim{The mixture representation~\eqref{eq:mixture} expresses \(F\) as
an average of indicators of the systems \(A_{\rs\omega}\), whose
thresholds are shifted at random by
\(h\log(\rs{\omega_\alpha}/\rs{\omega_0})\), where
\(\rs{\omega_0},(\rs{\omega_\alpha})_\alpha\) are independent unit
exponentials. To verify it, observe that, conditional on
\(\rs{\omega_0}\), the probability that \(x\in A_{\rs\omega}\)} is
\(e^{-\rs{\omega_0}T(x)}\), and its expectation is \((1+T(x))^{-1}\).
Equivalently, \skl{the variables \(-\log\omega_\beta\)} are independent standard
Gumbel variables, so
\[
 F_{A,h}(x)=\Pp\!\left\{\sks{-h\log\omega_0}\ge
 \max_\alpha(\nu_\alpha\cdot x-a_\alpha-\sks{h\log\omega_\alpha})\right\}.
\]
\rev{Here the baseline \trim{category} has deterministic score zero, and the face
displacements \(\sks{h\log(\omega_\alpha/\omega_0)}\) have logistic marginals and are
dependent through \(\sks{\omega_0}\); independent logistic displacements would give
a different approximation. McFadden~\cite[p.~1001]{McFadden1989} uses
the same random-utility representation, and Gao and
Zhang~\cite[Lemma~14]{GaoZhang} state it for their softmax weights.} \revnote{Citation of the unavailable
Gumbel note removed.}

\trim{The chain-rule formula sketched in
Subsection~\ref{sec:outline-smoothing} reads, for vectors}
\(z_1,\ldots,z_q\),
\begin{equation}
 \begin{split}
 D^qF(x)[z_1,\ldots,z_q]
 ={}&h^{-q}\sum_{\rs{\mathcal P}\in\rs{\Pi_q}}(-1)^{|\rs{\mathcal P}|}|\rs{\mathcal P}|!
 \sum_{\alpha_1,\ldots,\alpha_{|\rs{\mathcal P}|}}
 F(x)\prod_{b=1}^{|\rs{\mathcal P}|}\pi_{\alpha_b}(x)\\
 &\hspace{3em}\times
 \prod_{b=1}^{|\rs{\mathcal P}|}\prod_{a\in\rs{J_b}}
 (\nu_{\alpha_b}\cdot z_a),
 \end{split}
 \label{appx-eq:partition-derivative}
\end{equation}
where \(\rs{\Pi_q}\) is the set of partitions \(\rs{\mathcal P}\) of \(\{1,\ldots,q\}\),
and \(\rs{J_1},\ldots,\rs{J_{|\mathcal P|}}\) are the blocks in any fixed ordering.
Ordered partitions can be encoded by a permutation and cuts between
successive elements, so
\(\sum_{\rs{\mathcal P}\in\rs{\Pi_q}}|\rs{\mathcal P}|!\le q!2^{q-1}\).
In particular,
\begin{equation}
 |D^qF(x)[z_1,\ldots,z_q]|
 \le C^q q!h^{-q}\prod_{a=1}^q
 \max_\alpha|\nu_\alpha\cdot z_a|.
 \label{appx-eq:all-orders}
\end{equation}
\rev{Gao and Zhang~\cite[Lemma~13]{GaoZhang} obtain
\eqref{appx-eq:partition-derivative} and~\eqref{appx-eq:all-orders}
from the joint cumulants of a categorical variable.}

The uniform bound \eqref{appx-eq:all-orders} does not yet record that the
derivatives are small away from the boundary. \trim{For the ordered
coordinate tuples \(I\) and the derivatives \(\partial_I\) of
Subsection~\ref{sec:outline-smoothing}, the} same positive majorant,
using \(\norm{\nu_\alpha}_1=1\), gives
\begin{equation}
 \sum_I\sup_{\norm z_\infty\le\rs\varrho}|\partial_IF(x+z)|
 \le C_qh^{-q}F(x)(1-F(x)),\qquad 0\le\rs\varrho\le h.
 \label{appx-eq:local-envelope}
\end{equation}
To see the local stability, \pass{observe that} each \(w_\alpha\)
changes by a factor in
\([e^{-\rs\varrho/h},e^{\rs\varrho/h}]\), and \(1+T\) obeys the same comparison.
After summing over coordinate tuples and face indices, the terms with
\(k\) blocks are bounded by a constant times
\(h^{-q}T^k/(1+T)^{k+1}\le h^{-q}F(1-F)\).

\skl{In the next lemma, \(U\) has an independent component in the sense
that \(U=S+R\), where \(S\) is independent of \(R\) and \(S\) is either a
sum of independent centered vectors with summand bound \(2\delta\) and
equal coordinate variances at least \(1/2\), or a centered Gaussian vector
with equal coordinate variances at least \(1/2\).}

\begin{lemma}[Averaging the derivative envelope]\label{appx-lem:envelope}
Suppose a random vector \(U\) has an independent component that is
either a sum of independent centered vectors, with summand bound
\(2\delta\) and equal coordinate variances at least \(1/2\), or a
centered Gaussian vector with equal coordinate variances at least
\(1/2\).
Uniformly over thresholds, for \(h\ge4\delta\),
\begin{equation}
 \E F(U)(1-F(U))\le Ch\sqrt L.
 \label{appx-eq:small-envelope}
\end{equation}
Consequently, for fixed \(q\) and \(h\ge\max\{4\delta,\rs\varrho\}\),
\begin{equation}
 \E\sum_I\sup_{\norm z_\infty\le\rs\varrho}|\partial_IF(U+z)|
 \le C_q h^{1-q}\sqrt L.
 \label{appx-eq:averaged-derivatives}
\end{equation}
\end{lemma}

\begin{proof}
Shift every face threshold by the same amount \(s\), and write
\(F_s=(1+e^{-s/h}T)^{-1}\) for the resulting logistic approximation. Since
\(\partial_sF_s=F_s(1-F_s)/h\), and
\(F_s(1-F_s)\ge e^{-1}F(1-F)\) for \(|s|\le h\),
\[
 F(1-F)\le\frac e2(F_h-F_{-h}).
\]
In the mixture representation~\trim{\eqref{eq:mixture}},
\(F_h-F_{-h}\) is the average of
\(\ind_{A_{\rs\omega}^h}-\ind_{A_{\rs\omega}^{-h}}\), with shifts interpreted in the
\trim{face projections \(\nu_\alpha\cdot x\)}. Conditional on the exponentials and the
independent shift, Lemma~\rev{\ref{lem:concentration}} bounds the expectation
by \(C(h+\delta)\sqrt L\)\pass{, since} the base has equal coordinate
variances
at least \(1/2\) and summand bound at most \(2\delta\), or is
Gaussian. Since \(h\ge4\delta\), this proves
\eqref{appx-eq:small-envelope}. Combining it with
\eqref{appx-eq:local-envelope} proves \eqref{appx-eq:averaged-derivatives}.
\end{proof}

\subsection{The error at the endpoints}

\rev{The mixture representation gives an endpoint bound for every
covariance matrix, and a second argument, which uses the correlation
gap, can be sharper. We choose the better of the two after fixing the
smoothing scale.}

\begin{lemma}[Unrestricted endpoint]\label{appx-lem:endpoint-free}
Under \eqref{eq:setup}, for \(h\ge\delta\),
\begin{equation}
 \E|F_{A,h}(W)-\ind_A(W)|
 +\E|F_{A,h}(Z)-\ind_A(Z)|
 \le ChL^{3/2}.
 \label{appx-eq:endpoint-free}
\end{equation}
\end{lemma}

\begin{proof}
In \trim{\eqref{eq:mixture}}, let
\(R_h=h\max_\alpha|\log(\rs{\omega_\alpha}/\rs{\omega_0})|\).
Then \(A^{-R_h}\subseteq A_{\rs\omega}\subseteq A^{R_h}\), and
\(\E R_h\le Ch\log(em)\le ChL\).
The last assertion follows by integrating the exponential tails of
\(-\log\rs{\omega_\alpha}\), and the faster upper tails of \(\log\rs{\omega_\alpha}\).
For fixed exponentials, the indicators can disagree only in
\(A^{R_h}\setminus A^{-R_h}\). Jensen's inequality moves the absolute
value inside the exponential average. Conditional on the exponentials,
Lemma~\rev{\ref{lem:concentration}} bounds the errors at \(W\) and \(Z\)
by \(C(R_h+\delta)\sqrt L\) and \(CR_h\sqrt L\), respectively.
Averaging the radius gives \(C(hL+\delta)\sqrt L\) and
\(ChL\sqrt L\).
Use \(h\ge\delta\) and \(L\ge1\).
\end{proof}

\begin{lemma}[Curvature endpoint]\label{appx-lem:endpoint-gap}
Under \eqref{eq:setup}, if \(c>0\) and \(h\ge\delta\), then
\begin{equation}
 \E|F_{A,h}(W)-\ind_A(W)|
 +\E|F_{A,h}(Z)-\ind_A(Z)|
 \le \frac{Ch\sqrt L}{c}.
 \label{appx-eq:endpoint-gap}
\end{equation}
\end{lemma}

\begin{proof}
\rev{We first bound the error of \(F\) pointwise by the curvature
\(\chi\) of \(\Psi\), and then bound the expected curvature by means of
the correlation gap.} Direct differentiation gives
\begin{equation}
 D^2\Psi=\frac1h\left(
 \sum_\alpha\pi_\alpha\nu_\alpha\nu_\alpha^\top-\rs{\bar\nu}\rs{\bar\nu}^\top\right),
 \qquad \norm{\rs{\bar\nu}}_1\le\theta,\qquad
 \chi=\theta-\norm{\rs{\bar\nu}}_2^2.
 \label{appx-eq:hessian}
\end{equation}
Distinct normals \trs{\(\pm e_j\)} have nonpositive inner products.
Thus, with \(\pi_* =\max_\alpha\pi_\alpha\),
\(\chi\ge\theta(1-\theta)\) and
\(\chi\ge\theta(1-\pi_*)\).
On \(A\), every \(w_\alpha\le1\), so \(\pi_*\le1/2\).
Off \(A\) \pass{we have} \(T>1\), so \(\theta>1/2\). On \(A\), the second
lower bound gives \(\chi\ge\theta/2\); off \(A\), the first gives
\(\chi\ge F/2\). The indicator error is \(\theta\) and \(F\)
in the two cases, respectively, so
\begin{equation}
 |F-\ind_A|\le2\chi.
 \label{appx-eq:curvature-error}
\end{equation}
To bring the covariance into this estimate, observe that
\[
 \sum_{j\ne k}|\partial_{jk}\Psi|
 =h^{-1}(\norm{\rs{\bar\nu}}_1^2-\norm{\rs{\bar\nu}}_2^2)
 \le\tr D^2\Psi,
\]
and since \(\Sigma_{jj}=1\) and \(|\Sigma_{jk}|\le1-c\) for
\(j\ne k\), the off-diagonal contribution can remove at most
\((1-c)\tr D^2\Psi\). Therefore
\begin{equation}
 \tr(\Sigma D^2\Psi)\ge c\,\tr D^2\Psi.
 \label{appx-eq:gap-contraction}
\end{equation}

\rev{To compare the Hessian at nearby points, include \trim{the index \(0\), with}
\(\nu_0=0\) \trim{and \(\pi_0=F\)}; the positive representation}
\[
 D^2\Psi=\frac1{2h}\sum_{\alpha,\beta\in\trs{\{0,1,\ldots,m\}}}
 \pi_\alpha\pi_\beta
 (\nu_\alpha-\nu_\beta)(\nu_\alpha-\nu_\beta)^\top
\]
implies
\(e^{-4\rs\varrho/h}D^2\Psi(x)\preceq D^2\Psi(x+z)
\preceq e^{4\rs\varrho/h}D^2\Psi(x)\) when \(\norm z_\infty\le\rs\varrho\)\rev{;
here \(M\preceq N\) means that \(N-M\) is positive semidefinite}.
Set \(W_i=W-\xi_i\). Centering and independence let us subtract
\(\nabla\Psi(W_i)\) before integrating the Hessian along the segment
from \(W_i\) to \(W\). First compare this Hessian with
\(D^2\Psi(W_i)\). Independence then replaces
\(\xi_i\xi_i^\top\) by \(\E\xi_i\xi_i^\top\). A second comparison
replaces \(D^2\Psi(W_i)\) by \(D^2\Psi(W)\), and summing the \rs{covariances of the
summands} gives \(\Sigma\). Each lower comparison contributes a factor
\(e^{-4\delta/h}\), so
\[
 \begin{split}
 \E[W\cdot\nabla\Psi(W)]
 &=\sum_i\E\int_0^1
 \xi_i^\top D^2\Psi(W_i+s\xi_i)\xi_i\dd s\\
 &\ge e^{-8\delta/h}\E\tr(\Sigma D^2\Psi(W)).
 \end{split}
\]
Since \(\norm{\nabla\Psi}_1\le1\), the left side is at most
\(\E\norm W_\infty\), which Bernstein's inequality bounds by
\(C(\sqrt L+\delta L)\). When \(\delta\sqrt L\le1\), this is
\(C\sqrt L\), and \(h\ge\delta\) bounds \(e^{8\delta/h}\)
absolutely. Combining with \eqref{appx-eq:gap-contraction} gives
\(\E\chi(W)\le Ch\sqrt L/c\).
If \(\delta\sqrt L>1\), the same conclusion follows from
\(\chi\le1\) and \(h\ge\delta\).
For \(Z\), Gaussian integration by parts gives
\[
 \E[Z\cdot\nabla\Psi(Z)]
 =\E\tr(\Sigma D^2\Psi(Z)),\qquad
 \E\norm Z_\infty\le C\sqrt L.
\]
\rev{This identity holds also for singular \(\Sigma\), as one sees by
restricting to its support.} Use \eqref{appx-eq:curvature-error} to finish.
\end{proof}


\begin{revblock}
\begin{proof}[Proof of Lemma~\ref{lem:endpoint}]
Lemma~\ref{appx-lem:endpoint-free} bounds the endpoint error by
\(ChL^{3/2}=Ch\sqrt L\cdot L\), and, when \(c>0\),
Lemma~\ref{appx-lem:endpoint-gap} bounds it by \(Ch\sqrt L\cdot c^{-1}\).
The smaller of the two is \(Ch\sqrt L\,\kappa\), which
is~\eqref{eq:endpoint-error}. Since
\(|\Pp(W\in A)-\Pp(Z\in A)|\) is at most
\(|\E F_{A,h}(W)-\E F_{A,h}(Z)|\) plus the two endpoint errors,
\eqref{eq:endpoint-min} follows.
\end{proof}
\end{revblock}


% Source: covariance_support (rewritten for readability, 25 October 2026)
\begin{revblock}
\section{\texorpdfstring{\trs{Comparison along the path: proof of Proposition~\ref{an:prop:derivative}}}{Comparison along the path: proof of Proposition 1}}\label{appx-sec:support}

This appendix proves Proposition~\ref{an:prop:derivative} along the
outline of Subsections~\ref{sec:outline-early}
and~\ref{sec:outline-late}\trim{, with the derivative formulas and the
averaged envelope of Appendix~\ref{appx-sec:smoothing}; Appendix~\ref{sec:main-proof}
integrates the resulting bounds}. Its six subsections proceed in the following
order.
\begin{enumerate}
\item[\trs{\ref{appx-an:sec:coords}.}] \trim{We justify the support coordinates of
Subsection~\ref{sec:outline-late} and define concentrated bases.}
\item[\trs{\ref{appx-an:sec:path}.}] We write down the path and an exact formula for
\(P_A'(\lambda)\): a cubic term at the full path plus three
fourth-order corrections.
\item[\trs{\ref{appx-an:sec:early}.}] We prove the early bounds.
\item[\trs{\ref{appx-an:sec:restoration}.}] We prove the late bounds, taking for granted a matrix
estimate for the logistic weights. The main \trim{step} is an exact expansion
of the weights at translated points, which brings every derivative back
to a common base.
\item[\trs{\ref{appx-an:sec:matrix}.}] We prove the matrix estimate, taking for granted a Gaussian
surface estimate.
\item[\trs{\ref{appx-an:sec:surface}.}] We prove the surface estimate.
\end{enumerate}
Subsections~\trs{\ref{appx-an:sec:coords}--\ref{appx-an:sec:restoration}} give the structure of the argument, and
Subsections~\trs{\ref{appx-an:sec:matrix}} and~\trs{\ref{appx-an:sec:surface}} prove the two Gaussian estimates on which it
rests.

\subsection{Coordinates on the covariance support}
\label{appx-an:sec:coords}

\trim{Subsection~\ref{sec:outline-late} introduced the support
coordinates, namely the factor \(\Sigma=\Gamma\Gamma^\top\) with
\(\Gamma\) of full column rank \(r\), the summands \(\zeta_i^0\) with
\(\xi_i=\Gamma\zeta_i^0\), the standard Gaussian vector \(G\) on
\(\R^r\) with \(Z\stackrel d=\Gamma G\), the unit vectors
\(u_\alpha=\Gamma^\top\nu_\alpha\), and the face norm
\(\norm x_U=\max_\alpha|u_\alpha\cdot x|\). We justify them here.}
Every null direction \(z\) of \(\Sigma\) satisfies
\(\sum_i\E(z\cdot\xi_i)^2=z^\top\Sigma z=0\), so each summand \(\xi_i\) lies
almost surely in the range of~\(\Gamma\)\trim{, and
\(\sum_i\E\zeta_i^0(\zeta_i^0)^\top=I_r\) because \(\Gamma\) has full
column rank and
\(\Gamma\{\sum_i\E\zeta_i^0(\zeta_i^0)^\top\}\Gamma^\top
=\Sigma=\Gamma\Gamma^\top\)}. Since
\(\norm{u_\alpha}^2=\nu_\alpha^\top\Sigma\nu_\alpha\) is a diagonal entry
of~\(\Sigma\), every \(u_\alpha\) is a unit vector\trim{, and because}
\(u_\alpha\cdot\zeta_i^0=\nu_\alpha\cdot\xi_i\) is a signed
coordinate of an original summand,
\begin{equation}
 \norm{\zeta_i^0}_U\le\delta.
 \label{appx-eq:projected-bound}
\end{equation}
The face bound controls every derivative direction that is evaluated
against a face, and the identity
\trim{\(\sum_i\E\zeta_i^0(\zeta_i^0)^\top=I_r\)} lets us sum a pair of
summand directions without taking a trace\trim{; nothing else about the
size of the summands is used}.

\trim{As in Subsection~\ref{sec:outline-late}, \(F\), \(w_\alpha\), and
\(\pi_\alpha\) denote the restricted functions on \(\R^r\), so that
\(F(x)=F_{A,h}(\Gamma x)\) and \(w_\alpha=e^{t_\alpha}\) with the face
score \(t_\alpha(x)=(u_\alpha\cdot x-a_\alpha)/h\). All} the
derivative formulas of Appendix~\ref{appx-sec:smoothing} hold with
\(u_\alpha\) in place of~\(\nu_\alpha\). In particular, a derivative of
\(F\) in directions \(z_1,\ldots,z_q\) is the derivative of \(F_{A,h}\)
in the directions \(\Gamma z_1,\ldots,\Gamma z_q\).

\paragraph{Concentrated bases.}
\trim{We now make precise the notion of a concentrated base from
Subsection~\ref{sec:outline-late}, a random translation \(V_0\) of
\(\R^r\) over which several estimates below average.} For \(a\in\R^m\), where \(m\le2p\) is the number of faces, write
\[
 Q(a)=\{x:u_\alpha\cdot x\le a_\alpha,\ \alpha\le m\},
 \qquad Q(a)^t=Q(a+t\mathbf1).
\]
\trim{Thresholds for which \(Q(a)\) is empty} are allowed. We call \(V_0\) a
\emph{concentrated base} if, for an absolute constant \(C_0\),
\begin{equation}
 \Pp\{V_0\in Q(a)^t\setminus Q(a)^{-t}\}
 \le C_0(t+\delta)\sqrt L,
 \qquad a\in\R^m,\quad t>0.
 \label{appx-an:concentration}
\end{equation}
The property passes to subsystems, by sending unused thresholds to
infinity, and to deterministic translates of~\(V_0\). Under \(\Gamma\),
the event in~\eqref{appx-an:concentration} is a boundary \skl{layer} of a system
of coordinate inequalities. Lemma~\ref{lem:concentration}, applied
to the \trim{vector of face projections} as in Appendix~\ref{appx-sec:smoothing}, will
therefore show that every base used below is concentrated.
\trim{Subsections~\ref{appx-an:sec:restoration}--\ref{appx-an:sec:surface}
split the path as \(U=V_0+\sigma G\), with \(G\) standard Gaussian on
\(\R^r\) and independent of~\(V_0\). Every Gaussian computation there is
made with \(V_0\) fixed, and the average over \(V_0\) then gives the
factor \(h\sqrt L\) in~\eqref{appx-an:weight-concentration} below and
the factor \(\sigma\) in the surface estimate of
Lemma~\ref{appx-an:lem:surface}. The proof of the late bounds in
Subsection~\ref{appx-an:sec:restoration} shows that the path supplies
such a base.}

Over a concentrated base the logistic weights are small on average, in
the sense that, with \(\theta=1-F\) as in
Appendix~\ref{appx-sec:smoothing} and \(h\ge\delta\),
\begin{equation}
 \sup_z\E F(V_0+z)\theta(V_0+z)\le Ch\sqrt L.
 \label{appx-an:weight-concentration}
\end{equation}
The proof of Lemma~\ref{appx-lem:envelope} uses only the concentration
of its base and therefore applies here. It rests on the inequality
\(F(1-F)\le\frac e2(F_h-F_{-h})\) for the common threshold shifts
\(\pm h\), and on the fact that \(F_h-F_{-h}\) is an average of
indicators of \skl{boundary layers} of width~\(2h\), each of probability at most
\(C_0(h+\delta)\sqrt L\). The same bound holds for the function
\(e^{t_j}/(1+e^{t_j})^2\) of a single inequality, since it is \(F(1-F)\)
for the subsystem consisting of that face.

\subsection{The path and its derivative}
\label{appx-an:sec:path}

\trim{In support coordinates, the path~\eqref{eq:path} of
Subsection~\ref{sec:outline-path}, with the same \(\gamma\), \(\sks{\varpi}\), and
\(\sks{\mathbb B}_{i,\lambda}\), reads}
\begin{equation}
 Y_i=\gamma \sks{\mathbb B}_{i,\lambda}\zeta_i^0,\qquad V_\lambda=\sum_iY_i,
 \qquad U_\lambda=V_\lambda+\sqrt\lambda\,G,
 \qquad \trs{0\le\lambda\le1},
 \label{appx-an:path}
\end{equation}
where the selection variables \(\sks{\mathbb B}_{i,\lambda}\), the summands, and \(G\) are
independent. Only the law of \(U_\lambda\) at each \(\lambda\) is used,
so no coupling across \(\lambda\) is needed. Since \(\sks{\varpi}\gamma^2=1-\lambda\),
the covariance of \(U_\lambda\) is \(I_r\) for every \(\lambda\), and
\(\Gamma U_\lambda\) has the law of the path~\eqref{eq:path}; hence
\(P_A(\lambda)=\E F(U_\lambda)\). \trim{Since} the cubic coefficient
\(\sks{\varpi}\gamma^3=1-\lambda^2\) has zero derivative at \(\lambda=0\)\trim{, as
arranged in Subsection~\ref{sec:outline-path},} the cubic part of the
early bound~\eqref{an:early} vanishes linearly in~\(\lambda\).

To write the moments of the summands explicitly, let \(H_i=U_\lambda-Y_i\) be the
path with summand~\(i\) deleted, and let \(\zeta_i\) and
\(\widetilde\zeta_i\) be two independent copies of \(\zeta_i^0\),
independent of the path and of~\(G\). Thus \(\E\zeta_i^{\otimes3}\) is
the third-moment tensor of summand~\(i\).

The next lemma splits \(P_A'(\lambda)\) into a cubic term at the full
path and three fourth-order corrections.

\begin{lemma}[Exact fourth-order identity]\label{appx-an:lem:identity}
For \(0<\lambda<1\),
\begin{equation}
 P_A'(\lambda)
 =-\frac\lambda3\sum_i\E D^3F(U_\lambda)
                    [\zeta_i,\zeta_i,\zeta_i]
       +\sum_i(\mathcal A_i+\mathcal B_i+\mathcal C_i),
 \label{appx-an:fourth-identity}
\end{equation}
where
\begin{align}
 \mathcal A_i
 &=\int_0^1 k_\lambda(s)
       \E D^4F(H_i+s\gamma\zeta_i)
                 [\zeta_i,\zeta_i,\zeta_i,\zeta_i]\dd s,
 \label{appx-an:A}\\
 k_\lambda(s)
 &=\frac\gamma6(1-s)^2\{(3-\lambda)s-2\lambda\},
 \label{appx-an:kernel}\\
 \mathcal B_i
 &=\frac{1-\lambda}{2}\int_0^1(1-s)
       \E D^4F(H_i+s\gamma\widetilde\zeta_i)
          [\zeta_i,\zeta_i,\widetilde\zeta_i,\widetilde\zeta_i]
       \dd s,
 \label{appx-an:B}\\
 \mathcal C_i
 &=\frac{\lambda(1-\lambda)}{3\gamma}\int_0^1
       \E D^4F(H_i+s\gamma\widetilde\zeta_i)
          [\zeta_i,\zeta_i,\zeta_i,\widetilde\zeta_i]\dd s.
 \label{appx-an:C}
\end{align}
\end{lemma}
\lin{In~\eqref{appx-an:fourth-identity}, the term \(\mathcal A_i\) is
the Taylor remainder of the \(i\)th summand; the term \(\mathcal B_i\) arises
because the Gaussian heat term acts at the full path while the quadratic
part of the summand acts at the leave-one-out base; and the term
\(\mathcal C_i\) moves the cubic part of the summand from \(H_i\) back to
\(U_\lambda\), which places all the cubic tensors at one base, where
their aggregate can cancel.}
\rev{Steps~1 and~2 of the proof follow
Aronow and Lopatto's proof of their Lemma~4.1~\cite[eqs.~(20)
and~(21)]{AL}, with the same fourth-order kernel. They leave the cubic
term at the leave-one-out base; Step~3 moves it to the full path, which
gives the identity the form of Gao and Zhang's derivative
identity~\cite[Lemma~22, eq.~(15)]{GaoZhang} for their path.}

\begin{proof}
Write \(M_i=\E\zeta_i\zeta_i^\top\) and
\(\sks{M_i^{(3)}}=\E\zeta_i^{\otimes3}\). All differentiations below are
legitimate for fixed \(h>0\), since the number of summands is finite, the summands are
bounded, and \(F\) and its derivatives are bounded.

\emph{Step 1: the summand operators.} For a smooth function \(f\), the
\(i\)th summand acts through
\[
 \ps{\mathcal S_i}f(x)=(1-\sks{\varpi})f(x)+\sks{\varpi}\E f(x+\gamma\zeta_i).
\]
Since \(\gamma'=1\), differentiation in \(\lambda\) gives
\[
 \ps{\dot{\mathcal S}_i}f(x)=\sks{\varpi}'\{\E f(x+\gamma\zeta_i)-f(x)\}
                  +\sks{\varpi}\E Df(x+\gamma\zeta_i)[\zeta_i].
\]
Expand the difference by Taylor's formula through degree three, and the
derivative through degree two. Centering removes the linear terms. Since
\((\sks{\varpi}\gamma^2)'=-1\) and \((\sks{\varpi}\gamma^3)'=-2\lambda\), the result is
\[
 \ps{\dot{\mathcal S}_i}f
 =-\tfrac12 M_i:D^2f
   -\tfrac\lambda3 \sks{M_i^{(3)}}:D^3f
   +\int_0^1k_\lambda(s)
        \E D^4f(\,\cdot+s\gamma\zeta_i)[\zeta_i^4]\dd s,
\]
where \([\zeta_i^4]\) means four repeated arguments. The fourth-order
kernel is exactly
\[
 \frac{\sks{\varpi}'\gamma^4}{6}(1-s)^3+\frac{\sks{\varpi}\gamma^3}{2}(1-s)^2
 =\frac\gamma6(1-s)^2\{(3-\lambda)s-2\lambda\},
\]
because \(\sks{\varpi}'=-(3-\lambda)/\gamma^3\) and \(\sks{\varpi}\gamma^3=(1-\lambda)\gamma\).
The derivative of \(P_A\) collects \(\ps{\dot{\mathcal S}_i}\) applied with the other
summands and the Gaussian vector held fixed, that is, at the base~\(H_i\).
The fourth-order part of this expansion gives \(\mathcal A_i\).

\emph{Step 2: the quadratic terms.} The Gaussian heat derivative
contributes \(\tfrac12\sum_iM_i:\E D^2F(U_\lambda)\), because
\(\sum_iM_i=I_r\). Combined with the quadratic terms of the summands, this leaves
\[
 \frac12\sum_i M_i:
             \E\{D^2F(U_\lambda)-D^2F(H_i)\}.
\]
Conditional on \(H_i\), the deleted summand \(Y_i\) is centered and
independent of~\(H_i\). Its second-order Taylor formula, with
\(\sks{\varpi}\gamma^2=1-\lambda\), identifies this expression with
\(\sum_i\mathcal B_i\).

\emph{Step 3: moving the cubic term.} Finally,
\[
 -\frac\lambda3 \sks{M_i^{(3)}}:\E D^3F(H_i)
 =-\frac\lambda3 \sks{M_i^{(3)}}:\E D^3F(U_\lambda)
   +\frac\lambda3 \sks{M_i^{(3)}}:
          \E\{D^3F(U_\lambda)-D^3F(H_i)\}.
\]
A first-order integral expansion of the last difference is exactly
\(\mathcal C_i\), because \(\sks{\varpi}\gamma=(1-\lambda)/\gamma\). The first term
on the right is the cubic term of~\eqref{appx-an:fourth-identity}. This
proves the identity.
\end{proof}

Every base that occurs in \(\mathcal A_i\), \(\mathcal B_i\), and
\(\mathcal C_i\) is a translate of the full path:
\begin{equation}
 H_i+s\gamma\zeta_i=U_\lambda+v_i,\qquad v_i=s\gamma\zeta_i-Y_i,
 \qquad \norm{v_i}_U\le4\delta,
 \label{appx-an:translates}
\end{equation}
and similarly with \(\widetilde\zeta_i\) in place of~\(\zeta_i\). The
shift bound uses \(s\gamma\le2\) and \(\norm{Y_i}_U\le2\delta\).

\subsection{The early bounds}
\label{appx-an:sec:early}

Near \(\lambda=0\) the path is almost the bounded sum itself, so for the
early bounds \eqref{an:early} and~\eqref{eq:remainder-early} we estimate
the derivatives in~\eqref{appx-an:fourth-identity} directly, by their
size near the boundary of the rectangle. \rev{The argument follows the
rest of Aronow and Lopatto's proof of their Lemma~4.1~\cite{AL}. Step~2
below corresponds to their moment bounds~(22) and~(23), and Step~3 to
their local supremum over shifts of size~\(4\delta\); they reach a
common law through an independent copy of the deleted summand, where we use
the translates~\eqref{appx-an:translates}. They bound the local suprema
by Gaussian smoothing~\cite[Corollary~3.2]{AL}, which requires
\(\lambda_{\min}(\Sigma)>0\), and Steps~1 and~3 use the logistic
envelope of Lemma~\ref{appx-lem:envelope} instead.}

\emph{Step 1: the path is concentrated.} We first show that
\begin{equation}
 \E F(U_\lambda)\theta(U_\lambda)\le Ch\sqrt L,
 \qquad 0\le\lambda\le1.
 \label{appx-an:path-weight}
\end{equation}
For \(\lambda\le1/2\), the bounded part \(V_\lambda\) is independent of
the Gaussian part, has every face variance \(1-\lambda\ge1/2\), and has
summands bounded by \(2\delta\) in the face norm. For \(\lambda\ge1/2\), the
Gaussian part has every face variance \(\lambda\ge1/2\). In both cases
Lemmas~\ref{lem:concentration} and~\ref{appx-lem:envelope}
give~\eqref{appx-an:path-weight}.

\emph{Step 2: the moment coefficients.} For this step we return to the
original coordinates. Put \(\xi_i=\Gamma\zeta_i\) and
\(\widetilde\xi_i=\Gamma\widetilde\zeta_i\). These \trim{copies} are
independent of the full path \(\Gamma U_\lambda\), have maximum norm at
most \(\delta\), and satisfy \(\sum_i\E\xi_{ij}^2=1\) for every
coordinate~\(j\). For every coordinate triple, Cauchy--Schwarz gives
\[
 \sum_i\E|\xi_{ij}\xi_{ik}\xi_{i\ell}|
 \le\delta\sum_i(\E\xi_{ij}^2\E\xi_{ik}^2)^{1/2}
 \le\delta.
\]
In the same way, every fourth-order coefficient that occurs in
\eqref{appx-an:A}, \eqref{appx-an:B}, or~\eqref{appx-an:C}, after
absolute values are taken and the sum over \(i\) is formed, is at most
\(\delta^2\); bound two of the four factors by \(\delta\) each, and
apply the same Cauchy--Schwarz argument to the remaining pair.

\emph{Step 3: the local envelope.} By~\eqref{appx-an:translates}, every
base in the fourth-order terms is the full path translated by a shift
whose original-coordinate image \(\Gamma v_i=s\gamma\xi_i-\Gamma Y_i\)
has maximum norm at most \(4\delta\le h\). Take the supremum over all
such shifts in the local derivative envelope~\eqref{appx-eq:local-envelope}.
Since the \trim{copies that give the} directions are independent of the full path, the moment
coefficients of Step~2 separate from the derivatives. The cubic term
of~\eqref{appx-an:fourth-identity} is therefore at most
\(C\lambda\delta h^{-3}\E F\theta(U_\lambda)\), and the three
fourth-order terms together are at most
\(C\delta^2h^{-4}\E F\theta(U_\lambda)\); the scalar kernels have absolute
integrals on \([0,1]\) bounded uniformly in~\(\lambda\). Substituting
\eqref{appx-an:path-weight} gives
\[
 |R_{4,A,h}(\lambda)|\le\frac{C\delta^2\sqrt L}{h^3},
 \qquad
 |P_A'(\lambda)|\le C\sqrt L\left(\frac{\lambda\delta}{h^2}
                         +\frac{\delta^2}{h^3}\right),
\]
which are \eqref{eq:remainder-early} and~\eqref{an:early}.

\subsection{The late bounds}
\label{appx-an:sec:restoration}

Once the path contains Gaussian variance \(\sigma^2\) of order at least
\(\delta^2L\), two derivatives can be moved onto the Gaussian density.
Lemma~\ref{appx-an:lem:restoration} below does this for a sum over the summands
of derivatives taken at translated bases, and the late bounds then
follow by a suitable choice of base.

Lemma~\ref{appx-an:lem:restoration} rests on the matrix estimate proved
in Subsection~\ref{appx-an:sec:matrix}, which we state now. Let
\(V_0\) be a concentrated base, let \(G\sim N(0,I_r)\) be independent of
\(V_0\), let \(\delta\sqrt L\le\sigma\le1\), and put \(U=V_0+\sigma G\).
Lemma~\ref{appx-an:lem:matrix} asserts that, for \(h\ge\delta\) and every
integer \(k\ge1\),
\begin{equation}
 \E_{V_0}\sum_{\alpha_1,\ldots,\alpha_k}
 \left\|\E_G\left[(GG^\top-I_r)F(U)
                \prod_{a=1}^k\pi_{\alpha_a}(U)\right]\right\|_*
 \le C_1^{k+1}hL^{3/2},
 \label{appx-an:matrix-bound}
\end{equation}
where \(C_1\ge1\) is an absolute constant\trim{ and the} sum runs over
all \(m^k\) ordered tuples of faces.
\skl{In the next lemma, the random vectors \(Z_i\) and \(Z_{i,a}\) are
auxiliary arrays in \(\R^r\), unrelated to the Gaussian target~\(Z\).}

\begin{lemma}[Paired derivatives at translated bases]
\label{appx-an:lem:restoration}
Let \(\mathcal H\) be a \(\sigma\)-field, let \(V_0\) be an
\(\mathcal H\)-measurable concentrated base, let \(G\sim N(0,I_r)\) be
independent of \(\mathcal H\), let \(\delta\sqrt L\le\sigma\le1\), and put
\(U=V_0+\sigma G\). For \(s\in\{1,2\}\), let
\((Z_i,Z_{i,1},\ldots,Z_{i,s},v_i)_i\) be random arrays which, conditional
on \(\mathcal H\), are independent of~\(G\). They may depend on
\(\mathcal H\), and on one another. Suppose that, almost surely,
\begin{equation}
 \sum_i\E[Z_iZ_i^\top\mid\mathcal H]\preceq I_r,
 \qquad \norm{Z_{i,a}}_U\le\delta,
 \qquad \norm{v_i}_U\le b\delta.
 \label{appx-an:conditional-budget}
\end{equation}
There is an absolute constant \(\mathsf A_b\), depending only on \(b\),
such that \(h\ge\mathsf A_b\delta\) implies
\begin{equation}
 \left|\sum_i\E D^{s+2}F(U+v_i)
         [Z_i,Z_i,Z_{i,1},\ldots,Z_{i,s}]\right|
 \le\frac{C_b}{\sigma^2}\,
                 \delta^s h^{1-s}L^{3/2}.
 \label{appx-an:restoration-bound}
\end{equation}
\end{lemma}

\begin{proof}
We remove the repeated pair of directions by two Gaussian integrations
by parts, rewrite the translated weights in terms of the weights at the
common base \(U\), sum the coefficients of the summands \trim{using the
covariance bound in}~\trs{\eqref{appx-an:conditional-budget}},
and finally apply~\eqref{appx-an:matrix-bound}.

\emph{Step 1: integration by parts.} Condition on \(\mathcal H\) and on
the arrays. A derivative of \(F(V_0+\sigma g+v_i)\) in a direction \(z\)
is \(\sigma^{-1}\) times the derivative in \(g\), so two Gaussian
integrations by parts give
\[
 \begin{split}
 \E_G D^{s+2}F(U+v_i)[Z_i,Z_i,Z_{i,1},\ldots,Z_{i,s}]
 &=\sigma^{-2}\E_G\Bigl[
    \tr\{(GG^\top-I_r)Z_iZ_i^\top\}\\
 &\hspace{6em}\times D^sF(U+v_i)[Z_{i,1},\ldots,Z_{i,s}]\Bigr].
 \end{split}
\]
By the derivative formula~\eqref{appx-eq:partition-derivative}, each
partition \(\mathcal P=\{J_1,\ldots,J_k\}\) of \(\{1,\ldots,s\}\)
contributes to \(D^sF(x)[Z_{i,1},\ldots,Z_{i,s}]\) the term
\[
 h^{-s}(-1)^kk!\sum_{\alpha_1,\ldots,\alpha_k}
 F(x)\prod_{a=1}^k\pi_{\alpha_a}(x)
 \prod_{a=1}^k\prod_{\iota\in J_a}(u_{\alpha_a}\cdot Z_{i,\iota}).
\]

\emph{Step 2: translated weights.} \trim{We use the translation
formula~\eqref{eq:weight-translation} of
Subsection~\ref{sec:outline-late}.} For a shift \(v\), put
\[
 d_j(v)=e^{u_j\cdot v/h}-1\trs{,
 \quad\text{so that}\quad} z_v(x)=\sum_j\pi_j(x)d_j(v),
 \qquad \epsilon=e^{b\delta/h}-1.
\]
If \(\norm v_U\le b\delta\), then \(|d_j(v)|\le\epsilon\) and
\(|z_v(x)|\le\epsilon\). When \(\epsilon<1\), the \trim{factor
\(\{1+z_v\}^{-k-1}\) in~\eqref{eq:weight-translation}} has the absolutely
convergent expansion
\[
 \{1+z_v\}^{-k-1}
 =\sum_{\sks{t}\ge0}(-1)^{\sks{t}}\binom{k+\sks{t}}{\sks{t}}z_v^{\sks{t}},
 \qquad
 z_v^{\sks{t}}=\sum_{\beta_1,\ldots,\beta_{\sks{t}}}
                \prod_{q=1}^{\sks{t}}\pi_{\beta_q}(x)d_{\beta_q}(v).
\]
Thus a term of order \(\sks{t}\) carries the weight
\(F(U)\prod_a\pi_{\alpha_a}(U)\prod_q\pi_{\beta_q}(U)\), with \(k+\sks{t}\)
factors \(\pi\), all evaluated at the common base~\(U\).

\emph{Step 3: the sum over \(i\).} For fixed face tuples \(\alpha\) and
\(\beta\), the summands enter only through the coefficient matrix
\[
 \Xi_{\alpha,\beta}=\sum_i\E\left[
  Z_iZ_i^\top
  \prod_{a=1}^k\prod_{\iota\in J_a}(u_{\alpha_a}\cdot Z_{i,\iota})\,
  e^{\sum_a u_{\alpha_a}\cdot v_i/h}
  \prod_{q=1}^{\sks{t}} d_{\beta_q}(v_i)
       \,\middle|\,\mathcal H\right].
\]
For every deterministic vector \(z\), the bounds
in~\eqref{appx-an:conditional-budget} give, pointwise in~\(\mathcal H\),
\[
 |z^\top\Xi_{\alpha,\beta}z|
 \le\delta^se^{kb\delta/h}\epsilon^{\sks{t}}
               \sum_i\E[(z\cdot Z_i)^2\mid\mathcal H]
 \le\delta^se^{kb\delta/h}\epsilon^{\sks{t}}\norm{z}^2.
\]
The matrix is symmetric, so its operator norm obeys the same bound, for
every realization of~\(\mathcal H\). This step uses the covariance
\trim{bound} in place of a trace, so no factor~\(r\) appears.

\emph{Step 4: the matrix estimate.} By the duality between the
operator norm and the nuclear norm,
\(|\tr(\Xi M)|\le\norm{\Xi}_{\mathrm{op}}\norm{M}_*\). Apply it to each
pair of tuples, average over \(\mathcal H\), and
use~\eqref{appx-an:matrix-bound} with \(k+\sks{t}\) weights. Only the
averaged matrix estimate uses the concentration of~\(V_0\), so no
concentration conditional on \(\mathcal H\) is needed. The resulting
series is at most
\[
 e^{kb\delta/h}C_1^{k+1}
 \sum_{\sks{t}\ge0}\binom{k+\sks{t}}{\sks{t}}(C_1\epsilon)^{\sks{t}}
 =e^{kb\delta/h}C_1^{k+1}(1-C_1\epsilon)^{-k-1}
 \le C_b^{k+1},
\]
provided \(h\ge\mathsf A_b\delta\), which makes \(C_1\epsilon\le1/2\).
Since \(s\le2\), there are boundedly many partitions, with bounded
coefficients. Restoring the factor \(\sigma^{-2}h^{-s}\)
proves~\eqref{appx-an:restoration-bound}.

The interchanges of sums and expectations can first be justified with a
finite bound that depends on~\(r\). The absolute face sums at order
\(\sks{t}\) are at most \(\epsilon^{\sks{t}}\), and their binomial series sums to
\((1-\epsilon)^{-k-1}\). Moreover \(\sum_i\E\norm{Z_i}^2\le r\), so the
\trim{factor \(\tr\{(GG^\top-I_r)Z_iZ_i^\top\}\)} is integrable after summing over the arrays. This finite
bound serves only the justification and does not enter the estimate.
\end{proof}

\begin{proof}[Proof of the late bounds \eqref{an:late}
and~\eqref{an:remainder}]
Let \(d_0\delta^2L\le\lambda<1\), with \(d_0\ge2\).
Lemma~\ref{appx-an:lem:restoration} needs a concentrated base and an
independent Gaussian component for the integration by parts. For
\(\lambda\) near~\(1\) the bounded part has small variance and no longer
provides the concentration, so in that range we split the Gaussian
component in two.
\begin{itemize}
\item If \(\lambda\le1/2\), take \(V_0=V_\lambda\) and
\(\sigma^2=\lambda\), and let \(\mathcal H\) be generated by the entire
path array. The base is concentrated by Lemma~\ref{lem:concentration},
since its summands are bounded by \(2\delta\) in the face norm and every
face variance is \(1-\lambda\ge1/2\).
\item If \(\lambda\ge1/2\), write \(\sqrt\lambda\,G\) as the sum of two
independent Gaussian vectors with covariance \((\lambda/2)I_r\). Put one
of them into \(V_0\) and into \(\mathcal H\), and use the other as
\(\sigma G\), so that \(\sigma^2=\lambda/2\). The base is concentrated
by the Gaussian case of Lemma~\ref{lem:concentration}\trim{, which is
Nazarov's inequality}, since its Gaussian
component has every face variance at least \(1/4\).
\end{itemize}
In both cases \(\sigma^2\asymp\lambda\), and \(\lambda\ge d_0\delta^2L\)
ensures \(\sigma\ge\delta\sqrt L\).

For the cubic term of~\eqref{appx-an:fourth-identity}, apply
Lemma~\ref{appx-an:lem:restoration} with \(s=1\),
\(Z_i=Z_{i,1}=\zeta_i\), and \(v_i=0\). Because the copies are \trim{independent of \(\mathcal H\)},
\(\sum_i\E[\zeta_i\zeta_i^\top\mid\mathcal H]=I_r\). Consequently
\begin{equation}
 \left|\frac\lambda3\sum_i\E D^3F(U_\lambda)
                    [\zeta_i,\zeta_i,\zeta_i]\right|
 \le\frac{C\lambda}{\sigma^2}\delta L^{3/2}
 \le C\delta L^{3/2}.
 \label{appx-an:cubic-bound}
\end{equation}
For each fourth-order term, write its base as a translate of the full path
by~\eqref{appx-an:translates}, and apply the lemma with \(s=2\),
\(Z_i=\zeta_i\), and \(b=4\). The remaining two directions are
\((\zeta_i,\zeta_i)\) in \(\mathcal A_i\),
\((\widetilde\zeta_i,\widetilde\zeta_i)\) in \(\mathcal B_i\), and
\((\zeta_i,\widetilde\zeta_i)\) in \(\mathcal C_i\). Each has face norm
at most \(\delta\), and the same covariance \trim{bound} applies. Every summed
integrand is therefore at most \(C\delta^2L^{3/2}/(\lambda h)\) in
absolute value. The scalar kernels of \(\mathcal A_i\) and
\(\mathcal B_i\) have bounded absolute integrals, and the coefficient of
\(\mathcal C_i\) is bounded. This proves~\eqref{an:remainder}, and
together with~\eqref{appx-an:cubic-bound} it proves~\eqref{an:late}.
\end{proof}

Finally, the path is continuous in law at its endpoints: \(\Gamma U_0\)
has the law of \(W\), and \(U_1=G\) is standard Gaussian on the support.
Since \(F\) is bounded and continuous, \(P_A(0)=\E F_{A,h}(W)\) and
\(P_A(1)=\E F_{A,h}(Z)\). The early bound~\eqref{an:early} is integrable
on the whole interval, which justifies integrating the derivative up to
both endpoints. This completes the proof of
Proposition~\ref{an:prop:derivative}, apart from the matrix
estimate~\eqref{appx-an:matrix-bound}.

\subsection{A matrix estimate for the logistic weights}
\label{appx-an:sec:matrix}

To prove~\eqref{appx-an:matrix-bound}, we use, as in
Subsection~\ref{sec:outline-late}, that the Gaussian laws tilted by the
log-concave weights have covariance at most \(I_r\),
and that the nuclear norm of each matrix is then bounded by two scalar
quantities, the squared length of the tilted mean and a \trim{trace
term, which may have either sign}.
We control the means relative to a single face
(Lemma~\ref{appx-an:lem:mean}), and the \trim{trace term} by a Gaussian
surface estimate (Lemma~\ref{appx-an:lem:surface}). The argument is the
soft-weight analogue of the convex-cell bound of Gao and
Zhang~\cite[Lemma~12]{GaoZhang}, in which the weights are indicators of
convex cells.

Throughout, \(V_0\) is a concentrated base, \(G\sim N(0,I_r)\) is
independent of \(V_0\), \(\delta\sqrt L\le\sigma\le1\), and
\(U=V_0+\sigma G\). \trim{We use the face scores \(t_j\) of
Subsection~\ref{appx-an:sec:coords}, so that the denominator of \(F\) is
\(N=1+\sum_je^{t_j}\).}

A weight can move the Gaussian mean both normally and tangentially to a
face, and Lemma~\ref{appx-an:lem:mean} bounds the normal component by a
second moment and the tangential component by the loss of mass relative
to a weight that depends on the normal coordinate alone. \rev{It refines
the entropy inequality \(\norm{\mu_w}^2\le2\log(1/p_w)\), which Gao and
Zhang use in the proof of their Lemma~12~\cite{GaoZhang}, and which the
first step of the proof below gives in every direction when
\(b\equiv1\).}

\begin{lemma}[Gaussian mean relative to a single face]
\label{appx-an:lem:mean}
Let \(u\) be a unit vector, and let \(b(g)=\beta(u\cdot g)\in[0,1]\) have
positive Gaussian integral \(p_b\). If \(0\le w\le b\),
\(p_w=\E w(G)>0\), and \(\mu_w=\E[Gw(G)]/p_w\), then
\begin{equation}
 p_w\norm{\mu_w}^2
 \le \E[w(G)(u\cdot G)^2]
        +2p_w\log(p_b/p_w).
 \label{appx-an:mean-entropy}
\end{equation}
\end{lemma}

\begin{proof}
Let \(t\perp u\). Jensen's inequality, the bound \(w\le b\), and the
independence of the normal and tangential components of \(G\) give
\[
 p_we^{t\cdot\mu_w}
 \le\E[w(G)e^{t\cdot G}]
 \le\E[b(G)e^{t\cdot G}]
 =p_be^{\norm{t}^2/2}.
\]
Take \(t\) to be the projection of \(\mu_w\) onto \(u^\perp\); then the
tangential part of \(\mu_w\) satisfies
\(p_w\norm t^2\le2p_w\log(p_b/p_w)\). The normal part satisfies
\(p_w|u\cdot\mu_w|^2\le\E[w(G)(u\cdot G)^2]\) by Cauchy--Schwarz. Adding
the two estimates proves the lemma.
\end{proof}

\begin{lemma}[Matrix estimate for logistic weights]
\label{appx-an:lem:matrix}
If \(h\ge\delta\), then~\eqref{appx-an:matrix-bound} holds for every
integer \(k\ge1\).
\end{lemma}

\begin{proof}
Condition on \(V_0\), and for a tuple \(\alpha=(\alpha_1,\ldots,\alpha_k)\)
write
\[
 w_\alpha(g)=F(V_0+\sigma g)
                  \prod_{a=1}^k\pi_{\alpha_a}(V_0+\sigma g),
 \qquad p_\alpha=\E_Gw_\alpha(G),
 \qquad \ps{\bar w}(g)=\sum_\alpha w_\alpha(g)=F(U)\theta(U)^k.
\]

\emph{Step 1: the tilt and the trace reduction.} The logarithm of
\(w_\alpha\) is a linear function of \(g\) minus
\((k+1)\log N(V_0+\sigma g)\), and \(\log N\) is convex, so \(w_\alpha\)
is log-concave. Let \(\mu_\alpha\) and \(\sks{\Sigma_\alpha}\) be the mean and
covariance of the Gaussian law tilted by \(w_\alpha\). Pr\'ekopa's
theorem~\cite[Theorem~6]{Prekopa1973} states that integrating a
log-concave function of two vectors over one of them leaves a
log-concave function of the other. Applied to the jointly log-concave
function of \((z,g)\) under the integral, it shows that the function
\[
 z\longmapsto\int_{\R^r}e^{-\norm{g-z}^2/2}w_\alpha(g)\dd g
\]
is log-concave. The Hessian of its logarithm at zero is
\(\sks{\Sigma_\alpha}-I_r\), so \(\sks{\Sigma_\alpha}\preceq I_r\). \rev{For indicators of
convex cells, Gao and Zhang~\cite[Lemma~11]{GaoZhang} obtain the same
bound from the Brascamp--Lieb inequality\pass{, a variance bound for log-concave
laws}.} Now
\(\E_G[(GG^\top-I_r)w_\alpha(G)]=p_\alpha(\sks{\Sigma_\alpha}-I_r+\mu_\alpha\mu_\alpha^\top)\).
Since \(I_r-\sks{\Sigma_\alpha}\succeq0\), the nuclear norm of the right side is at
most \(p_\alpha\{\tr(I_r-\sks{\Sigma_\alpha})+\norm{\mu_\alpha}^2\}\), and its trace is
\(p_\alpha(\tr \sks{\Sigma_\alpha}-r+\norm{\mu_\alpha}^2)=\E_G[H_2(G)w_\alpha(G)]\),
where \(H_2(g)=\norm g^2-r\). Eliminating \(\tr \sks{\Sigma_\alpha}\) gives
\begin{equation}
 \left\|\E_G[(GG^\top-I_r)w_\alpha(G)]\right\|_*
 \le2p_\alpha\norm{\mu_\alpha}^2
             -\E_G[H_2(G)w_\alpha(G)].
 \label{appx-an:matrix-trace-reduction}
\end{equation}
\trim{The function \(H_2\) enters only as the trace of
\(gg^\top-I_r\). If the weights \(w_\alpha\) formed a partition of
unity, as the indicators of convex cells do in Gao and Zhang's
Lemma~12, the trace terms would add up to \(\E H_2(G)=0\). Here they add
up to \(\E_G[H_2(G)F(U)\theta(U)^k]\), which Step~4 bounds.}
\skl{We now} bound the sum over \(\alpha\) of each term on the right.

\emph{Step 2: the means.} Use the first face of each tuple as a
one-dimensional reference. For a tuple whose first entry is \(j\), put
\[
 b_j(g)=\frac{e^{t_j(V_0+\sigma g)}}
                   {\{1+e^{t_j(V_0+\sigma g)}\}^2},
 \qquad p_{b_j}=\E_Gb_j(G).
\]
This function depends only on \(u_j\cdot g\), and
\(w_\alpha\le F\pi_j\le b_j\) because \(N\ge1+e^{t_j}\). Since
\(|u_j\cdot G|\le\norm G_U\), Lemma~\ref{appx-an:lem:mean} gives
\begin{equation}
 \sum_\alpha p_\alpha\norm{\mu_\alpha}^2
 \le\E_G[\ps{\bar w}(G)\norm G_U^2]
       +2\sum_\alpha p_\alpha\log(p_{b_{\alpha_1}}/p_\alpha).
 \label{appx-an:entropy-sum}
\end{equation}
Conditionally on \(G\), the bound \(\ps{\bar w}\le F\theta\)
and~\eqref{appx-an:weight-concentration} show that the first term has
expectation at most \(Ch\sqrt L\,\E\norm G_U^2\le ChL^{3/2}\).

\emph{Step 3: the entropy term.} Since there are \(m^{k-1}\) tuples with
a given first face, the log-sum inequality
\(\sum_j\ps{y_j\log(z_j/y_j)\le y\log(z/y)}\), for positive
\(\ps{y_j,z_j}\) with sums \(\ps y\) and~\(\ps z\), lets us add their reference masses before
estimating the logarithm. Put
\[
 \bar p=\E_{V_0,G}\ps{\bar w}(G),\qquad
 \bar b=m^{k-1}\sum_j\E_{V_0,G}b_j(G).
\]
The log-sum inequality, followed by the joint concavity of
\((x,y)\mapsto x\log(y/x)\), gives
\[
 \E_{V_0}\sum_\alpha p_\alpha
          \log(p_{b_{\alpha_1}}/p_\alpha)
 \le\bar p\log(\bar b/\bar p).
\]
By~\eqref{appx-an:weight-concentration} and its version \trim{for a
single face}, applied
conditionally on \(G\), we have \(\bar p\le Ch\sqrt L\) and
\(\bar b\le Cm^kh\sqrt L\). Write \(\bar p=h\sqrt L\,x\), with
\(0<x\le C\). Then
\[
 \bar p\log(\bar b/\bar p)
 \le h\sqrt L\,x\log(Cm^k/x)
 \le C(k+1)hL^{3/2},
\]
since \(\log m\le CL\). The case \(\bar p=0\) follows by continuity.

\emph{Step 4: the \trim{trace term}.} Summing the weights first gives
\(\sum_\alpha w_\alpha=F\theta^k\)\trim{, before any absolute value is
taken}. Shift every threshold by \(s\), and write
\(F_s,\theta_s\) for the resulting functions. Then
\[
 F\theta^k=-\frac hk\left.\partial_s\theta_s^k\right|_{s=0},
\]
so the \trim{trace term} is a derivative in a common threshold shift. To turn
it into a surface integral, we represent powers of \(F_s\) as mixtures of
polyhedral indicators. For an integer \(\sks{t}\ge1\), let
\(T_{\sks{t}}\sim\operatorname{Gamma}(\sks{t},1)\) and let \(\omega_j\) be
independent unit exponentials, all independent of \(V_0\) and~\(G\). The
Laplace transform of the Gamma law gives
\begin{equation}
 F_s(x)^{\sks{t}}=
 \Pp\{u_j\cdot x\le a_j+s+h\log(\omega_j/T_{\sks{t}}),\ j\le m\}.
 \label{appx-an:gamma-mixture}
\end{equation}
Expand \(\theta_s^k=\sum_{\sks{t}=0}^k(-1)^{\sks{t}}\binom k{\sks{t}} F_s^{\sks{t}}\). The
absolute coefficients sum to \(2^k\), and the term \(\sks{t}=0\) has zero
derivative. The first-variation formula at the end of
Subsection~\ref{appx-an:sec:surface}, which expresses the derivative of a
Gaussian integral over a polyhedron, when all its faces move outward at
unit speed, as an integral over its boundary, the mixture~\eqref{appx-an:gamma-mixture},
and Lemma~\ref{appx-an:lem:surface} yield
\[
 \left|\E_{V_0,G}H_2(G)F(U)\theta(U)^k\right|
 \le\frac{2^kh}{k\sigma}
   \sup_Q\E_{V_0}|\surf((Q-V_0)/\sigma)|
 \le C2^khL^{3/2}.
\]
The common shift \(-h\log T_{\sks{t}}\) has an absolutely continuous law, so
the first-variation formula applies\trim{, although the polyhedron
may become lower dimensional at exceptional values of the shift}.

Combining Steps~1--4 proves~\eqref{appx-an:matrix-bound}.
\end{proof}

\subsection{An averaged Gaussian surface estimate}
\label{appx-an:sec:surface}

Step~4 of the proof of Lemma~\ref{appx-an:lem:matrix} used a bound for
\trim{the} Gaussian surface integrals \trim{\(\surf\) of
Subsection~\ref{sec:outline-late}} of randomly translated polyhedra,
which we now prove. \trim{We keep the notation \(H_2\), \(\varphi_r\),
\(\mathcal H^{r-1}\), facets, and ridges of that subsection, and write
\(\Phi_r\) for the standard Gaussian measure on \(\R^r\).}
\trim{The weight \(H_2(g)=\norm g^2-r\) comes from the trace term in
Step~1 of that proof, since \(H_2\) is the trace of \(gg^\top-I_r\). The
translation is the concentrated base \(V_0\), and \(\sigma\) is the scale
of the Gaussian part of the path. A bound uniform over translations,
such as the bound \(C\{\log(2m)\}^{3/2}\) of Gao and
Zhang~\cite[Proposition~4]{GaoZhang}, would cost a factor \(1/\sigma\)
in the late bounds, and averaging over \(V_0\) removes it
(Subsection~\ref{sec:why-rate}).} Call a
polyhedron \emph{proper} if it is neither empty nor the whole space, and
\emph{full dimensional} if it has an interior point. \trim{The integral
\(\surf(Q)\) is defined for} proper, full-dimensional \trim{polyhedra} \(Q\)
with normals among \(u_1,\ldots,u_m\), \trim{with} each facet counted
once\trim{, and we} set \(\surf(Q)=0\) when \(Q\) is empty or the whole
space. \lin{We call the polyhedra \(Q\) whose rescaled translates
\((Q-V_0)/\sigma\) are averaged target polyhedra.}

\begin{lemma}[Averaged scalar surface]\label{appx-an:lem:surface}
Let \(V_0\) be a concentrated base, and let \(0<\sigma\le1\). Uniformly
over target polyhedra \(Q\) with normals among \(u_1,\ldots,u_m\), empty
or full dimensional,
\begin{equation}
 \E\left|\surf\bigl((Q-V_0)/\sigma\bigr)\right|
 \le C\bigl(\sigma L^{3/2}+\delta L^2\bigr).
 \label{appx-an:surface-bound}
\end{equation}
In particular, if \(\sigma\ge\delta\sqrt L\), the right side is at most
\(C\sigma L^{3/2}\).
\end{lemma}

The lemma omits lower-dimensional targets, which occur only at
exceptional values of a common threshold shift; the shifts used in
Step~4 above have absolutely continuous laws, and the first-variation
formula below makes this precise.

\begin{geoblock}
\paragraph{Terms from the geometry of polyhedra.}
A \emph{polyhedron} is an intersection of finitely many closed half-spaces
\(\{x:u_j\cdot x\le a_j\}\), where the unit vectors \(u_j\) are the
\emph{outward normals} and the numbers \(a_j\) the \emph{offsets}. An
inequality is \emph{redundant} if deleting it leaves the set unchanged. Let
\(Q\) be a proper, full-dimensional polyhedron in \(\R^r\), described without
redundant inequalities. Each inequality then defines a \emph{facet}
\(F_j=Q\cap\{u_j\cdot x=a_j\}\), a piece of the boundary of dimension
\(r-1\) that lies in one hyperplane, and the boundary \(\partial Q\) is the
union of the facets. Two facets that meet in a set of dimension \(r-2\)
share a \emph{ridge}. In the plane the facets are the sides of a polygon and
the ridges its corners; in \(\R^3\) the facets are the faces of a solid and
the ridges its edges; on the line the facets are the endpoints of an
interval, and there are no ridges. The measure \(\mathcal H^{r-1}\) is area
on a facet and \(\mathcal H^{r-2}\) the corresponding measure on a ridge.
Every point outside \(Q\) has a unique \emph{nearest point} in \(Q\). The
\emph{normal cone} of \(Q\) at a boundary point \(x\) is the set of
directions \(v\) for which \(x\) is the nearest point of \(x+tv\) for every
\(t>0\). At a point lying on the facet \(j\) only, it is the ray spanned by
\(u_j\); at a point lying on the facets \(j\) and \(k\) only, it is the
sector between \(u_j\) and \(u_k\) in the plane that they span. The
\emph{outward tube of length \(\varepsilon\)} over a facet or a ridge is
the set of points \(x+v\) with \(x\) in the facet or ridge, \(v\) in the
normal cone at \(x\), and \(\norm v<\varepsilon\).
\end{geoblock}

\begin{proof}
\begin{geoblock}
Condition on \(V_0\) and put \(Q'=(Q-V_0)/\sigma\), so that
\(V_0+\sigma G\in Q\) exactly when \(G\in Q'\). The set \(Q'\) is again a
polyhedron with normals among \(u_1,\ldots,u_m\), and it is empty or full
dimensional. If it is empty or the whole space, then \(\surf(Q')=0\) and
there is nothing to bound, so assume that \(Q'\) is proper and full
dimensional. Since \(\surf(Q')\) depends on the set \(Q'\) and not on its
description, we may delete every inequality that is redundant or repeats
another one. After this, each inequality defines exactly one facet, different
facets have different outward normals, and two facets that share a ridge
have normals that are neither equal nor opposite, because facets in parallel
hyperplanes do not meet in a ridge of a full-dimensional polyhedron; the
angle between such normals lies in \((0,\pi)\).

The proof first writes
\(\surf(Q')\) as an explicit sum of facet terms and ridge terms (Step~1),
then bounds each term by the Gaussian mass of a thin outward tube
(Step~2), and finally sums the tubes over dyadic sizes and averages over
\(V_0\) (Step~3).
\end{geoblock}
\rev{Steps~1 and~2 follow
Gao and Zhang's proof of their Proposition~4~\cite{GaoZhang}, in the
case where each matrix weight there is the identity. The divergence
identity of Step~1 combines their Lemmas~24, 25, and~27, and the tube
estimates of Step~2 adapt their Lemmas~28 and~29. Step~3, which averages
over the translation and gains the factor~\(\sigma\), has no counterpart
there.}

\emph{Step 1: a divergence identity.}
\begin{geoblock}
Index the facets of \(Q'\) by \(j\), and write \(F_j\) for the facet with
outward normal \(u_j\), lying in the hyperplane
\(\{g:u_j\cdot g=a'_j\}\). The offset \(a'_j\) is the signed distance from the
origin to this hyperplane, whose point nearest the origin is \(a'_ju_j\); it
is an offset of \(Q'\), not of \(Q\). Let
\(s_j=\int_{F_j}\varphi_r\dd\mathcal H^{r-1}\) be the Gaussian surface mass of
facet~\(j\). For a ridge \(E_{jk}\) shared by the facets \(j\) and \(k\), let
\(\alpha_{jk}=\arccos(u_j\cdot u_k)\in(0,\pi)\) be the angle between their
normals, let \(\rho_{jk}=\int_{E_{jk}}\varphi_r\dd\mathcal H^{r-2}\), and let
\(z_{jk}\) be the point of \(\operatorname{span}(u_j,u_k)\) with
\(u_j\cdot z_{jk}=a'_j\) and \(u_k\cdot z_{jk}=a'_k\). Every point \(g\) of the
ridge satisfies these two equations, so \(z_{jk}\) is the orthogonal
projection of every point of the ridge onto \(\operatorname{span}(u_j,u_k)\),
and it is the point of the intersection of the two hyperplanes nearest the
origin. Then
\begin{equation}
 \surf(Q')=
 \sum_j\bigl((a'_j)^2-1\bigr)s_j
 -\sum_{\{j,k\}}(d_{k\mid j}+d_{j\mid k})\rho_{jk},
 \label{appx-an:surface-identity}
\end{equation}
where the second sum runs over the ridges (the corners, when \(r=2\)), and
\[
 d_{k\mid j}=
 \frac{a'_k-(u_j\cdot u_k)a'_j}{\sin\alpha_{jk}},
 \qquad
 d_{k\mid j}+d_{j\mid k}=
 2\sin(\alpha_{jk}/2)\,
    \frac{u_j+u_k}{\norm{u_j+u_k}}\cdot z_{jk}.
\]

The identity is the divergence theorem, applied facet by facet. A point of
facet \(j\) is \(g=a'_ju_j+y\) with \(y\perp u_j\), and then
\(\varphi_r(g)=\varphi_1(a'_j)\varphi_{r-1}(y)\) and
\(H_2(g)=\{(a'_j)^2-1\}+\{\norm y^2-(r-1)\}\). The first part contributes
\(((a'_j)^2-1)s_j\). For the second part,
\[
 (\norm{y}^2-(r-1))\varphi_{r-1}(y)
 =-\operatorname{div}_{u_j^\perp}
                  \{y\varphi_{r-1}(y)\},
\]
the divergence being taken within the hyperplane \(u_j^\perp\), which is
parallel to the facet. The facet is itself a polyhedron in its hyperplane,
so the divergence theorem (Gauss--Green formula) turns the integral over
\(F_j\) into an integral over its boundary, the union of its ridges. Across
the ridge shared with facet \(k\), the unit vector in \(u_j^\perp\) that points
out of \(F_j\) is \(\nu_{k\mid j}=(u_k-(u_j\cdot u_k)u_j)/\sin\alpha_{jk}\), and
at every point \(g\) of the ridge
\(y\cdot\nu_{k\mid j}=g\cdot\nu_{k\mid j}=d_{k\mid j}\), because
\(g\cdot u_j=a'_j\) and \(g\cdot u_k=a'_k\). Facet \(j\) therefore contributes
\(-d_{k\mid j}\rho_{jk}\) at this ridge, and the same ridge contributes
\(-d_{j\mid k}\rho_{jk}\) when it is seen from facet~\(k\); adding the
contributions gives~\eqref{appx-an:surface-identity}. The second form of
\(d_{k\mid j}+d_{j\mid k}\) follows from \(1-\cos\alpha=2\sin^2(\alpha/2)\),
\(\sin\alpha=2\sin(\alpha/2)\cos(\alpha/2)\), \(\norm{u_j+u_k}=2\cos(\alpha/2)\),
and \((u_j+u_k)\cdot z_{jk}=a'_j+a'_k\).

Truncation justifies the calculation for unbounded facets, since polynomial
Gaussian terms vanish at infinity. When \(r=1\) there are no ridges, and the
identity is read as a sum over endpoints.

The two terms of a ridge must be paired before they are estimated. Each alone contains the factor
\(1/\sin\alpha_{jk}\), which is large when the two facets are nearly
parallel, while their sum is at most \(2\sin(\alpha_{jk}/2)\norm{z_{jk}}\le
\alpha_{jk}\norm{z_{jk}}\) in absolute value.
\end{geoblock}

\emph{Step 2: local tube estimates.} We bound each term
of~\eqref{appx-an:surface-identity} by the Gaussian mass of a short outward
tube. Fix \(R\ge2\) and put \(\varepsilon_R=(1+R)^{-1}\). Above a facet
with \(|a'_j|\le R\), the outward normal tube of length \(\varepsilon_R\)
has Gaussian mass \(q_j\) satisfying
\[
 q_j\ge c_*\varepsilon_Rs_j,
 \qquad (1+(a'_j)^2)s_j\le CR^3q_j,
\]
for a small absolute constant \(c_*\), because the density ratio at a
facet point displaced by \(tu_j\) is \(\exp(-ta'_j-t^2/2)\), which is
bounded below on the tube. For a ridge with \(\norm{z_{jk}}\le R\), use
displacements of norm less than \(\varepsilon_R\) in its two-dimensional
outward normal cone, the set of directions in which the ridge points are
the nearest points of \(Q'\); it is a sector of angle \(\alpha_{jk}\). For a
displacement \(y\) in that cone, \(x\cdot y=z_{jk}\cdot y\), so the same
density-ratio argument gives
\[
 q_{jk}\ge c_*\alpha_{jk}\varepsilon_R^2\rho_{jk}.
\]
Since \(|d_{k\mid j}+d_{j\mid k}|\le C\alpha_{jk}R\), it follows that
\(|d_{k\mid j}+d_{j\mid k}|\rho_{jk}\le CR^3q_{jk}\).

The tubes are disjoint up to null sets, because their points have
distinct metric projections, or nearest points, in \(Q'\). This uses two facts about
polyhedra: every ridge lies in exactly two facets, and the normal cone at
a relative-interior point of the ridge shared by facets \(j\) and \(k\) is
the cone spanned by \(u_j\) and~\(u_k\). All the tubes lie in
\((Q')^{\varepsilon_R}\setminus Q'\). These assertions concern the actual
faces of the polyhedron, so they \pass{hold also} when the original
threshold
system contains redundant inequalities.

\emph{Step 3: dyadic summation.} Divide the facet and ridge offsets into
dyadic shells. Let \(R_{\sks{\iota}}=2^{\sks{\iota}+1}\) for \(\sks{\iota}\ge0\). The zeroth shell
contains the facet offsets of absolute value at most \(2\) and the ridge
offsets of norm at most \(2\); shell \(\sks{\iota}\ge1\) contains the offsets in
\((R_{\sks{\iota}}/2,R_{\sks{\iota}}]\). Let \(\mathcal O_{\sks{\iota}}(V_0)\) be the union of the
associated tubes, all of radius \((1+R_{\sks{\iota}})^{-1}\). Then
\begin{equation}
 |\surf(Q')|\le C\sum_{\sks{\iota}\ge0}R_{\sks{\iota}}^3\,
                         \Phi_r(\mathcal O_{\sks{\iota}}(V_0)).
 \label{appx-an:dyadic-tubes}
\end{equation}
To average over \(V_0\), we record two consequences of the event
\(G\in\mathcal O_{\sks{\iota}}(V_0)\). Let \(\Pi_{jk}\) be the orthogonal
projection onto \(\operatorname{span}(u_j,u_k)\) for linearly independent
pairs, and put
\[
 \mathcal M(G)=\max\Bigl\{\max_j|u_j\cdot G|,
   \max_{j,k:\,u_j,u_k\text{ independent}}\norm{\Pi_{jk}G}\Bigr\},
\]
omitting the second maximum when it is empty. First, if
\(G\in\mathcal O_{\sks{\iota}}(V_0)\) and \(\sks{\iota}\ge1\), then
\(\mathcal M(G)\ge R_{\sks{\iota}}/2-1\), since the relevant facet coordinate or
ridge projection of \(G\) differs from the shell offset by at most one.
Second, for every \(\sks{\iota}\),
\[
 G\in\mathcal O_{\sks{\iota}}(V_0)
 \quad\Longrightarrow\quad
 V_0+\sigma G\in Q^{\sigma/(1+R_{\sks{\iota}})}\setminus Q.
\]
Condition on \(G\), apply~\eqref{appx-an:concentration} to the last event,
and retain the necessary event involving \(\mathcal M(G)\). This gives
\[
 \E\,\Phi_r(\mathcal O_{\sks{\iota}}(V_0))
 \le C(\sigma/R_{\sks{\iota}}+\delta)\sqrt L\,\beta_{\sks{\iota}},
\]
where \(\beta_0=1\) and \(\beta_{\sks{\iota}}=\Pp\{\mathcal M(G)\ge R_{\sks{\iota}}/2-1\}\)
for \(\sks{\iota}\ge1\). The maximum defining \(\mathcal M\) involves at most
\(m+m^2\) one- or two-dimensional Gaussian projections, so
\[
 \beta_{\sks{\iota}}\le\min\{1,Cm^2e^{-c_*R_{\sks{\iota}}^2}\},\qquad \sks{\iota}\ge1.
\]
Splitting the dyadic sum where \(R_{\sks{\iota}}\) is of order
\(\sqrt{\log(em)}\) yields, for \(q=2,3\),
\[
 \sum_{\sks{\iota}\ge0}R_{\sks{\iota}}^q\beta_{\sks{\iota}}
 \le C_q\{\log(em)\}^{q/2}\le C_qL^{q/2}.
\]
Taking expectations in~\eqref{appx-an:dyadic-tubes} and using Tonelli's
theorem proves~\eqref{appx-an:surface-bound}.
\end{proof}

\paragraph{First variation and common shifts.}
Step~4 of the proof of Lemma~\ref{appx-an:lem:matrix} differentiates in a
random common shift of all thresholds. The following facts justify this,
even when a polyhedron collapses at an exceptional shift. For fixed
offsets and \(V_0\), the function
\[
 s\longmapsto
 \E_G\left[H_2(G)\ind\left\{
   \max_j\{u_j\cdot(V_0+\sigma G)-a_j\}\le s\right\}\right]
\]
is absolutely continuous, because the inverse image of a Lebesgue-null
set under the maximum is contained in the union of its inverse images
under the individual face coordinates, each of which is a nondegenerate
Gaussian variable. The coarea formula, which writes the derivative in
\(s\) of an integral over the set where a Lipschitz function is at most
\(s\) as an integral over the level set where it equals \(s\), therefore
gives the derivative, for almost every \(s\), as
\(\sigma^{-1}\surf((Q(a+s\mathbf1)-V_0)/\sigma)\). To justify interchanges
before the dimension-free estimate is applied, a finite bound
\(Cm(r+1)/\sigma\) on the Lipschitz constant suffices: the integral of
\(|H_2|\) over a hyperplane at standard Gaussian coordinate \(t\) is at
most \((t^2+2r)\varphi_1(t)\le C(r+1)\), and there are at most \(m\)
facets. Independent absolutely continuous common shifts avoid the
exceptional levels and allow these first variations to be averaged. The
temporary dependence on \(r\) in this justification does not enter any
asserted bound.
\end{revblock}


% Source: concentration_proof
\section{\texorpdfstring{\trs{Proof of the anti-concentration inequality of Lemma~\ref{lem:concentration}}}{Proof of the anti-concentration inequality of Lemma 1}}\label{appx-app:concentration}

\trim{This appendix gives the two proofs of
Lemma~\ref{lem:concentration} described after its statement in
Subsection~\ref{sec:outline-smoothing}. Both use the quoted inequalities
as established results.}

\begingroup
\ifshowpolish\color{magenta}\fi
\subsection{Conditioning and Rademacher anti-concentration}
\label{appx-conc:conditioning}
\endgroup

We prove Lemma~\rev{\ref{lem:concentration}} by representing each summand
as a random choice between two independent copies. Conditional on these
copies, the sum becomes a translated Rademacher sum, to which we can
apply the anti-concentration inequality of O'Donnell, Servedio, and
Tan~\cite{OST} in the form recorded in~\cite[Lemma~J.4]{CCKK2022}.
This gives a short proof of the bounded-sum estimate stated
in~\cite[Lemma~2.1]{AL}.

First take \(v=1\) and omit the independent shift. Let
\(\zeta_i'\) be independent copies of the summands, and let
\(\varepsilon_i\) be independent uniform signs, independent of all
these vectors. Put
\[
 \rs{\bar\zeta_i}=\frac{\zeta_i+\zeta_i'}2,\qquad
 \rs{\Delta_i}=\frac{\zeta_i-\zeta_i'}2,\qquad
 \mathcal G=\sigma(\zeta_i,\zeta_i':1\le i\le n).
\]
Each \(\rs{\bar\zeta_i}+\varepsilon_i\rs{\Delta_i}\) selects one of the two copies, so
\begin{equation}
 S\stackrel d=\widetilde S:=\sum_i\rs{\bar\zeta_i}+\sum_i\varepsilon_i\rs{\Delta_i}.
 \label{appx-eq:pair-representation}
\end{equation}
Conditional on \(\mathcal G\), the first sum is a fixed translation,
\(\norm{\rs{\Delta_i}}_\infty\le\eta\), and the coordinate variances of the
second sum are \(V_j=\sum_i\rs\Delta_{ij}^2\). We need these variances to
stay away from zero. Independence and centering give
\[
 \E V_j=\frac12,\qquad
 0\le\rs\Delta_{ij}^2\le\eta^2,\qquad
 \sum_i\Var(\rs\Delta_{ij}^2)\le\frac{\eta^2}{2}.
\]
Bernstein's inequality and a union bound therefore imply, for an
absolute constant \(c_1>0\),
\begin{equation}
 \Pp(\mathcal E^c)\le p\exp(-c_1/\eta^2),
 \qquad \mathcal E=\{\min_{j\le p}V_j\ge1/4\}.
 \label{appx-eq:conditional-variance}
\end{equation}

On \(\mathcal E\), apply~\cite[Lemma~J.4]{CCKK2022} conditionally to
\(\sum_i\varepsilon_i\rs{\Delta_i}\). Encode the lower and upper coordinate
bounds by the \(2p\) coordinates of \((x,-x)\), and absorb the fixed
translation \(\sum_i\rs{\bar\zeta_i}\) into the thresholds. The outer and inner
rectangles correspond to thresholds separated by \(2\rs t\), so
\[
 \Pp(\widetilde S\in A^{\rs t}\setminus A^{-\rs t}\mid\mathcal G)
 \le C(\rs t+\eta)\sqrt L\qquad\text{on }\mathcal E.
\]
Infinite endpoints follow by taking limits. This argument also covers
inconsistent coordinate bounds, since the orthant inequality is uniform
in all thresholds. Averaging and using
\eqref{appx-eq:conditional-variance} gives
\begin{equation}
 \Pp(S\in A^{\rs t}\setminus A^{-\rs t})
 \le C(\rs t+\eta)\sqrt L+p\exp(-c_1/\eta^2).
 \label{appx-eq:concentration-pairs}
\end{equation}
If \(\eta\sqrt L\le\rs{c_2}\), with \(\rs{c_2}>0\) sufficiently small, then
\(L\le c_1/(2\eta^2)\), and the last term is at most
\(\exp(-c_1/(2\eta^2))\le C\eta\le C\eta\sqrt L\).
Otherwise the desired estimate follows from the bound by one.
This proves the bounded-sum assertion for \(v=1\). Dividing by
\(\sqrt v\) gives the general case, and conditioning on \(R\)
permits any independent shift.

For a centered Gaussian vector \(G\), Nazarov's \rev{anti-concentration}
inequality in
\cite[Lemma~A.1]{CCK2017}, also recorded in
\cite[Lemma~J.3]{CCKK2022}, applied to \((G,-G)/\sqrt v\) gives
\[
 \Pp(G\in A^{\rs t}\setminus A^{-\rs t})\le C\rs t\sqrt{L/v}.
\]
It requires only positive marginal variances and therefore permits
singular covariance. Conditioning on the independent shift and taking
the minimum with one completes the proof.


% Restored from the repository's original manuscript, with explicit attribution.
\begingroup
\ifshowpolish
\color{magenta}
\hypersetup{linkcolor=magenta,citecolor=magenta,urlcolor=magenta}
\fi
\subsection{Aronow--Lopatto's resampling proof}
\label{appx-conc:resampling}

\rev{This subsection gives the resampling proof of Aronow and
Lopatto~\cite[Lemma~2.1]{AL}, which compares two ways of measuring a
boundary layer.} Resampling a small
proportion of the \rs{summands} gives a point near a face a fixed positive chance
of leaving the rectangle, apart from a tail event. Kane's noise-sensitivity
bound controls the total probability of this change in membership.
Specifically, \cite[Corollary~4 in the arXiv version]{Kane} states that,
if \(f\) is the indicator of an intersection
of \(k\) affine halfspaces on a uniform Boolean cube\rev{, the set of
sign vectors \(\{-1,1\}^N\) with the uniform law}, and
\(\varepsilon^{\rs\vartheta}\) is obtained by independently flipping signs
with probability \(\rs\vartheta\in(0,1/2]\), then
\begin{equation}
 \Pp\{f(\varepsilon)\ne f(\varepsilon^{\rs\vartheta})\}
 \le C\sqrt{\rs\vartheta\log(ek)}.
 \label{appx-eq:kane}
\end{equation}

First normalize to \(v=1\) and omit the independent shift. Thus
\(S=\sum_i\zeta_i\) has independent centered \rs{summands},
\(\norm{\zeta_i}_\infty\le\eta\), and \(\Var(S_j)=1\) for every \(j\).
Independently replace each \rs{summand} by a fresh copy with probability
\(\rs\vartheta\), obtaining \(S^{\rs\vartheta}\).
To apply \eqref{appx-eq:kane}, generate two independent candidate copies
\(\zeta_i^+,\zeta_i^-\), and choose between them using a uniform sign.
Conditional on all candidate pairs, membership in a rectangle is an
intersection of at most \(2p\) affine halfspaces in these signs.
A sign flip selects the other independent candidate, so averaging
\eqref{appx-eq:kane} gives
\begin{equation}
 \Pp\{\ind_Q(S)\ne\ind_Q(S^{\rs\vartheta})\}
 \le C\sqrt{\rs\vartheta L}
 \label{appx-eq:resampling}
\end{equation}
for every rectangle \(Q\). Open or closed faces cause no difficulty on
a finite Boolean cube, since thresholds can be perturbed without changing
the Boolean function.

To show that points near a face can leave the rectangle, condition on
the original \rs{summands}. Use the equivalent representation
\(S^{\rs\vartheta}=\sum_i[(1-B_i)\zeta_i+B_i\zeta_i']\), where the
\(B_i\sim\operatorname{Bernoulli}(\rs\vartheta)\) and the fresh \rs{summands} are
independent of \(\rs{\mathcal H}=\sigma(\zeta_1,\ldots,\zeta_n)\).
For \(D=S^{\rs\vartheta}-S\), each coordinate obeys
\[
 D_j+\rs\vartheta S_j=\sum_iY_{ij},\qquad
 Y_{ij}=B_i(\zeta_{ij}'-\zeta_{ij})+\rs\vartheta\zeta_{ij}.
\]
The \(Y_{ij}\) are conditionally independent and centered,
\(|Y_{ij}|\le2\eta\), and
\[
 v_j:=\sum_i\E(Y_{ij}^2\mid\rs{\mathcal H})
 =\rs\vartheta+\rs\vartheta(1-\rs\vartheta)\sum_i\zeta_{ij}^2\ge\rs\vartheta.
\]
Also, \(\sum_i\E(|Y_{ij}|^3\mid\rs{\mathcal H})\le2\eta v_j\).
Let \(\Phi\) denote the standard normal distribution function.
The scalar Berry--Esseen inequality therefore yields, almost surely,
\begin{equation}
 \sup_x\left|
 \Pp\left\{\frac{D_j+\rs\vartheta S_j}{\sqrt{v_j}}\le x
 \,\middle|\,\rs{\mathcal H}\right\}-\Phi(x)\right|
 \le C\eta/\sqrt{\rs\vartheta}.
 \label{appx-eq:conditional-be}
\end{equation}

Choose the resampling variance large enough to cross a layer of width
\(d\), while keeping the Berry--Esseen error small. Let \(\partial_d Q\)
be the points in \(Q\) within distance \(d\) of a finite face. Put
\(\rs{\bar d}=d+\eta\) and \(\rs\vartheta=\rs{k_0}\rs{\bar d}^2\), with \(\rs{k_0}\) a large absolute
constant. We may assume \(\rs{\bar d}\sqrt L\) is sufficiently small, since
otherwise the desired estimate follows from the bound by one. This
assumption also ensures \(\rs\vartheta\le1/2\).
On
\(\{S\in\partial_d Q,\ \norm S_\infty\le1/(4\sqrt{\rs\vartheta})\}\),
select the first face within distance \(d\). This choice is
\(\rs{\mathcal H}\)-measurable, so \eqref{appx-eq:conditional-be} applies to
its coordinate and outward direction. Conditional on \(\rs{\mathcal H}\), the standardized displacement needed to cross that
face outward is at most
\[
 d/\sqrt{\rs\vartheta}+\sqrt{\rs\vartheta}|S_j|
 \le\rs{k_0}^{-1/2}+1/4.
\]
The error in \eqref{appx-eq:conditional-be} is at most \(C\rs{k_0}^{-1/2}\).
Choose \(\rs{k_0}\) large enough that the conditional outward-crossing
probability is bounded below by an absolute positive constant.
Crossing forces \(S^{\rs\vartheta}\notin Q\). Equation~\eqref{appx-eq:resampling}
then gives
\[
 \Pp(S\in\partial_d Q)
 \le C\rs{\bar d}\sqrt L
 +\Pp\{\norm S_\infty>1/(4\sqrt{\rs\vartheta})\}.
\]
Bernstein's inequality and a union bound imply
\[
 \Pp\{\norm S_\infty>1/(4\sqrt{\rs\vartheta})\}
 \le2\exp(L-c_1/\rs\vartheta).
\]
Here \(\eta/\sqrt{\rs\vartheta}\le\rs{k_0}^{-1/2}\) controls the bounded-summand
term in Bernstein's denominator. Decreasing the smallness threshold for
\(\rs{\bar d}\sqrt L\) ensures \(L\le c_1/(2\rs\vartheta)\), and therefore this
tail is at most \(C\rs{\bar d}\sqrt L\). We have proved
\begin{equation}
 \Pp(S\in\partial_d Q)\le C(d+\eta)\sqrt L.
 \label{appx-eq:inner-layer}
\end{equation}

If \(A^{\rs t}\) is empty, the desired bound is immediate. Otherwise every
point of \(A^{\rs t}\setminus A^{-\rs t}\) lies within distance \(2\rs t\)
of a finite face of \(A^{\rs t}\). If an interval of \(A^{-\rs t}\) is empty,
the corresponding interval of \(A^{\rs t}\) has length at most \(4\rs t\), so
the same assertion holds. Apply \eqref{appx-eq:inner-layer} with
\(Q=A^{\rs t}\) and \(d=2\rs t\). An independent shift is handled by
conditioning, and equal coordinate variances \(v\) by dividing by
\(\sqrt v\). This proves the bounded-sum part of
Lemma~\rev{\ref{lem:concentration}}.

For the Gaussian assertion, factor its covariance as
\(\Sigma=\sum_{k=1}^r \rs{f_k}\rs{f_k}^\top\), and approximate the Gaussian by
\[
 G_m=\frac1{\sqrt m}\sum_{k=1}^r\sum_{l=1}^m \rs{f_k}\varepsilon_{kl},
\]
where the \(\varepsilon_{kl}\) are independent uniform signs.
\pass{Each \(G_m\) has} covariance \(\Sigma\), while the maximum
summand norm tends to zero\pass{ as \(m\to\infty\)}. Every coordinate has variance \(v>0\), so each finite coordinate face
has Gaussian probability zero, even if \(\Sigma\) is singular. Each
fixed rectangle boundary layer is therefore a continuity set for the
Gaussian law\rev{, a set whose boundary has Gaussian probability zero, so
that convergence in distribution carries over to its probability}. The finite-dimensional CLT passes the bounded-sum estimate
to the limit. Conditioning again permits independent shifts.
\rev{This paragraph reproduces Aronow and Lopatto's proof of their
Corollary~2.2~\cite{AL}.}


\endgroup

% Source: subexponential_extension_proof
\section{\texorpdfstring{\trs{Sub-exponential coordinates: proof of Corollary~\ref{cor:subexponential}}}{Sub-exponential coordinates: proof of Corollary 2}}\label{app:subexponential}
\begin{proof}[Proof of Corollary~\ref{cor:subexponential}]
We clip each coordinate, apply the bounded theorem after marginal
standardization, and transfer the result back to \(W_n\). This requires
control of both the discarded sum and the Gaussian covariance change.
\trim{We drop the subscript \(n\) and write \(\ell\) and \(b\) for the
quantities \(\ell_n\) and \(b_n\) of the corollary.}

We first prove the bounds under \(\E e^{|X_{ij}|/K}\le2\) for every
coordinate. The stated norm condition gives this exponential-moment
bound at scale \(2K\): the defining infimum supplies an admissible
scale below \(2K\) for each coordinate, and increasing the scale
decreases its exponential moment. Replacing \(K\) by \(2K\) doubles
\(b\), while the positive logarithms in both asserted rates decrease.
Each rate therefore increases by at most a factor two.

The unrestricted bound follows from the bound by one, after adjusting its absolute
constant, when \(b\ge1/16\), so assume \(b<1/16\).

\emph{Step 1: Clip the \rs{summands} and bound the residual sum.}
\skl{With the clipping level \(4K\ell\),} define coordinatewise
\[
 C_T(x)=\operatorname{sgn}(x)\min\{|x|,\sks{4K\ell}\},\qquad
 Y_i=C_T(X_i)-\E C_T(X_i),\qquad \sks{D}_i=X_i-C_T(X_i).
\]
Here \(Y_i\) is the centered clipped \rs{summand} and \(\sks{D}_i\) is the discarded
part. Their normalized sums\skl{, marked by the subscripts \(T\) and \(D\),} satisfy
\[
 W_T=n^{-1/2}\sum_iY_i,\qquad
 \sks{W_D}=n^{-1/2}\sum_i(\sks{D}_i-\E \sks{D}_i),\qquad W_n=W_T+\sks{W_D}.
\]
The exponential-moment assumption gives, for \(t\ge0\),
\[
 \Pp(|\sks{D}_{ij}|>t)\le2e^{-\sks{4\ell}-t/K},\qquad
 \E|\sks{D}_{ij}|^k\le2k!K^ke^{-\sks{4\ell}}\quad(k\ge1).
\]
In particular, independence of the \rs{summands} implies
\begin{equation}
 \max_j\Var(\sks{W_{D,j}})\le4K^2e^{-\sks{4\ell}}.
 \label{tail:residual-moments}
\end{equation}
The pointwise bound \(e^{|\sks{D}_{ij}|/K}\le1+e^{-\sks{4\ell}}e^{|X_{ij}|/K}\)
gives \(\E e^{|\sks{D}_{ij}|/K}\le1+2e^{-\sks{4\ell}}\). Together with
\(\E|\sks{D}_{ij}|\le2Ke^{-\sks{4\ell}}\), independence and the triangle inequality
therefore yield
\begin{align*}
 \E\exp\!\left(\frac{\sqrt n\,|\sks{W_{D,j}}|}{K}\right)
 &\le\exp\!\left(\frac1K\sum_i\E|\sks{D}_{ij}|\right)
          \prod_i\E e^{|\sks{D}_{ij}|/K}\\
 &\le e^{2ne^{-\sks{4\ell}}}(1+2e^{-\sks{4\ell}})^n
 \le e^{4ne^{-\sks{4\ell}}}.
\end{align*}
For \(\sks{\bar L}=\log(4np)\), Markov's inequality gives
\[
 \Pp\!\left(\sqrt n\,|\sks{W_{D,j}}|\ge K[\sks{\bar L}+4ne^{-\sks{4\ell}}]\right)\le e^{-\sks{\bar L}}.
\]
Since \(n,p\ge1\),
\[
 4ne^{-\sks{4\ell}}=4e^{-4}n^{-3}\le\frac14\le\log4\le \sks{\bar L}.
\]
A coordinate union bound now yields
\begin{equation}
 \Pp(\|\sks{W_D}\|_\infty>\sks{\varrho_1})\le\frac1{4n}\le\frac1{2n},\qquad
 \sks{\varrho_1}=\frac{2K\sks{\bar L}}{\sqrt n}.
 \label{tail:sum-coupling}
\end{equation}
\rev{Since the union bound is applied to the residual sums, the clipping
level can depend only on \(\log n\).}

\emph{Step 2: Apply the bounded comparison.}
Set \(\tau=2Ke^{-\sks{2\ell}}\)\skl{, and write \(\operatorname{sd}(W_{T,j})\) for the
standard deviation of the \(j\)th coordinate of~\(W_T\)}.
Since each coordinate of \(W_n\) has variance one, the triangle
inequality in \(L^2\) gives
\[
 |\sks{\operatorname{sd}(W_{T,j})}-1|\le\tau,
 \qquad \frac\tau b=\frac{2\sqrt n\,e^{-2\ell}}\ell
 \le\frac2{e^2}.
\]
Hence \(\sks{\operatorname{sd}(W_{T,j})}\ge1/2\). Since \(\|Y_i\|_\infty\le\sks{8K\ell}\), the
independent centered \rs{summands} of the marginally standardized \(W_T\)
have envelope at most
\[
 \frac{\sks{8K\ell}}{\sqrt n\min_j\sks{\operatorname{sd}(W_{T,j})}}\le16b.
\]
Let \(Z_T\) be Gaussian with covariance \(\Cov(W_T)\).
Positive diagonal scaling preserves rectangle distance. Applying
Theorem~\ref{thm:unrestricted} with envelope \(16b\) gives
\begin{equation}
 \rh(W_T,Z_T)\le C\min\left\{1,
 bL^{3/2}\sqrt{1+\log_+\frac1{b\sqrt L}}\right\}.
 \label{tail:bounded-comparison}
\end{equation}
The factor 16 is absorbed into the absolute constant, since increasing
the envelope decreases the logarithmic factor.

\emph{Step 3: Restore the covariance.}
Construct a jointly centered Gaussian vector \((Z_T,\sks{Z_D})\) with the
block covariance of \((W_T,\sks{W_D})\). This covariance is positive
semidefinite, even when singular. Set \(Z_n=Z_T+\sks{Z_D}\) on this coupling.
Then \(Z_n\) has exactly the covariance of \(W_n\), while
\(\max_j\Var(\sks{Z_{D,j}})\le\tau^2\). Gaussian tails therefore give
\begin{equation}
 \Pp(\|\sks{Z_D}\|_\infty>\sks{\varrho_2})\le\frac1{2n},\qquad
 \sks{\varrho_2}=\sqrt8\,Ke^{-\sks{2\ell}}\sqrt{\sks{\bar L}}.
 \label{tail:Gaussian-coupling}
\end{equation}
The Gaussian components need not be independent.

For any rectangle \(A\), the two coupling bounds give
\begin{align*}
 \Pp(W_n\in A)
 &\le\Pp(W_T\in A^{\sks{\varrho_1}})+\frac1{2n}\\
 &\le\Pp(Z_T\in A^{\sks{\varrho_1}})+\rh(W_T,Z_T)+\frac1{2n}\\
 &\le\Pp(Z_n\in A^{\sks{\varrho_1}+\sks{\varrho_2}})+\rh(W_T,Z_T)+\frac1n.
\end{align*}
Inward expansions give the reverse bound
\(\Pp(W_n\in A)\ge\Pp(Z_n\in A^{-(\sks{\varrho_1}+\sks{\varrho_2})})
-\rh(W_T,Z_T)-1/n\). Applying Gaussian rectangle concentration at
the unit-marginal target \(Z_n\), from Lemma~\ref{lem:concentration}, yields
\begin{equation}
 \rh(W_n,Z_n)\le\rh(W_T,Z_T)
       +C(\sks{\varrho_1}+\sks{\varrho_2})\sqrt L+\frac1n.
 \label{tail:transfer}
\end{equation}
No independence between the two vectors within either coupling is
required for this sandwich argument.
The transfer error has order \(bL^{3/2}\): since
\(\sks{\bar L}\le2\ell L\) and \(K\ge1\),
\[
 (\sks{\varrho_1}+\sks{\varrho_2})\sqrt L+\frac1n
 \le CK\left[e^{-2\ell}\sqrt{L\sks{\bar L}}
                     +\frac{\sks{\bar L}\sqrt L}{\sqrt n}\right]+\frac1n
 \le CbL^{3/2}.
\]
Combine this with \eqref{tail:bounded-comparison} and the bound by one
to prove \eqref{eq:subexponential}.

\emph{Step 4: Preserve a fixed correlation gap.}
For \eqref{eq:subexponential-gap}, it suffices to consider
\(b<\min\{1/16,c_0/4\}\); otherwise its right-hand side is at
least one after increasing \(C_{c_0}\).
Under the Gaussian coupling, both \(Z_n\) and the standardized clipped
vector \((Z_{T,j}/\sks{\operatorname{sd}(W_{T,j})})_{j\le p}\) have unit coordinate variances, and
\[
 \|Z_{T,j}/\sks{\operatorname{sd}(W_{T,j})}-Z_{n,j}\|_{L^2}\le2\tau.
\]
Writing the difference of two covariance entries as a sum of two
products and applying Cauchy--Schwarz shows that the covariance matrices
differ entrywise by at most \(4\tau\). Since
\(4\tau\le(8/e^2)b<c_0/2\), the standardized clipped covariance
has correlation gap at least \(c_0/2\).
Theorem~\ref{thm:optimized} applied with envelope \(16b\) now gives
\[
 \rh(W_T,Z_T)\le C_{c_0}\min\left\{1,
 b\left[L^{3/2}+L\sqrt{1+\log_+\frac1{bL}}\right]\right\}.
\]
The transfer error in \eqref{tail:transfer} is bounded by the
\(CbL^{3/2}\) term. This proves \eqref{eq:subexponential-gap}.
\end{proof}

Coordinatewise odd clipping preserves central symmetry of each \rs{summand},
but need not preserve cancellation of third\pass{-moment} tensors across \rs{summands}.
Moreover the transfer cost above contains \(bL^{3/2}\), so the
zero-skewness refinement cannot simply be read off by deleting the
cubic term from the displayed sub-exponential bound.

\paragraph{Comparison with the locally nondegenerate result.}
The companion paper~\cite[Corollary~1]{Local} assumes the same
coordinatewise \(\psi_1\) condition, an averaged fourth-moment bound
\(\max_jn^{-1}\sum_i\E|X_{ij}|^4\le B_4^2\), and uniformly positive
eigenvalues for principal covariance submatrices of order at most four.
When these principal-submatrix eigenvalues are bounded below by \(a>0\), its bound is
\[
 \frac{C_a}{\sqrt n}\left[B_4L^{3/2}
                   +K\log(enp)\sqrt L\right],
\]
capped at one. For fixed \(K,B_4,a\) and \skl{\(p\) growing polynomially in~\(n\)},
this is \(O(n^{-1/2}(\log n)^{3/2})\), whereas
\eqref{eq:subexponential} is \(O(n^{-1/2}(\log n)^3)\).
The latter requires neither a fourth-moment hypothesis in addition to
the tail bound nor a local covariance assumption. If \(K\) varies,
the elementary fourth-moment consequence is \(B_4\le C K^2\);
the tail bound alone does not justify taking \(B_4=K\).


\end{document}

